% !TEX TS-program = xelatex
% !TEX encoding = UTF-8 Unicode

%下面四行用以创建hanout版本
%\documentclass[handout,cjkutf8,table]{beamer} % handout, slidestop 
%\usepackage{pgfpages}
%\pgfpagesuselayout{8 on 1}[a4paper,border shrink=5mm]
%\setbeameroption{show notes} % 显示note或只显示note(show only note)

\documentclass[handout,cjkutf8,table,aspectratio=1610]{beamer} % slidestop
%\mode* %mode*屏蔽所有不在frame中的内容，当需要生成讲稿模式的文件时，取消注释，并取消导言区关于article的三行注释
\input{../preamble.tex}
\usepackage{graphicx} % Allows including images
\usepackage{animate}
%--------------------------------------------------------------
%仅对本文件有效的导言
\newcounter{chp}\usecounter{chp}\setcounter{chp}{2} %设置一个章节计数器，
%\includeonlylecture{Lecture 4}
%\includeonlyframes{current} % 只显示带标签的页面
%\setbeameroption{show notes} % 显示note或只显示note(show only note)
%\renewcommand{\note}{a}
%\useoutertheme{shadow}
%%% Begin Document %%%
\begin{document}
%---------------------------------------------------------------------------------------------------
%% Set title
\title[自动控制原理]{自动控制原理}
\subtitle{第二章~~控制系统的数学模型}
\author[第二章~~控制系统的数学模型]{\vspace{-4em}\center{\includegraphics[scale=0.3]{../figs/JUSTlogo.png}}\\\vs{-3em}\large 电子信息学院\\ 张永韡~~博士~~副教授}\vspace{-4em}
\date[2021-2022学年春]{2021-2022学年春}
%---------------------------------------------------------------------------------------------------
%\title[自动控制原理]{\vspace{5em}《自动控制原理》讲稿}
%\author[第一章~~自动控制原理的一般概念]{\vspace{5em}\Large 第一章~~自动控制原理的一般概念}%文章模式
%---------------------------------------------------------------------------------------------------
\begin{frame}[plain]
	\titlepage
\end{frame}
%\maketitle
%--------------------------------------------------------------------------------------------
%:Lecture 1: 引言
\lecture{引言}{Lecture 1}
\begin{frame}{主要内容}
	\tableofcontents[hideallsubsections]
\end{frame}
%\note{
%\textbf{目标}
%了解建立系统模型的方法，使用微分方程建立典型环节的数学模型。\\
%
%电路系统；机械系统：1）位移系统，2）旋转系统；电枢控制直流电动机；位置控制系统；小结。

%\textbf{材料}
%
%1)RC网络$\rightarrow$R-L-C无源网络；2)弹簧-质量-阻尼器；3)旋转系统；
%4)电枢控制直流电机；
%5)位置控制系统。}
%
%\note{其他问题：\\
%迟交作业按最后成绩80\%打分，直到扣到及格为止。不交作业该次作业成绩记0分。\\
%作业抄袭按照作业成绩80\%打分，最低为及格。\\
%直接转发作业的很恶劣。\\
%以后为了区别作业次数，采用自控作业第X次+学号+姓名的格式\\
%如果上次布置作业没有通知最后交作业的期限，则延长一周作业时间，下次课之前交上\\
%公布交作业同学的名字，没交的同学自查是否没有发到，可能判别为垃圾邮件\\
%邮箱的问题(生活，工作)，邮箱姓名的问题-吴俊柯，有同学直接没有写姓名\\
%邮件格式，称呼，结束，标题zikongzuoye，附件名称\\
%文档格式，pdf保持格式统一。字体，字号，绘图\\
%最后，教学效果的反馈，可以使用生活邮箱\\
%几个亮点：温度，材料，底面积，接口吻合度等}

\section{引言}
\subsection{建模方法}
\begin{frame}{建模方法}
\begin{block}{机理建模（分析法）}
根据现有的物理、化学定律对系统各部分的动态进行分析和描述，并给出运动方程。
\end{block}
\begin{block}{电学：基尔霍夫定律}
\begin{description}
\item[电流定律] 所有进入某节点的电流的总和等于所有离开这节点的电流的总和。
\item[电压定律] 沿着闭合回路所有元件两端的电势差（电压）的代数和等于零。
\end{description}
\end{block}
\end{frame}

\begin{frame}{建模方法}
\begin{block}{力学：牛顿定律}
\begin{description}
\item[第一定律]存在某些参考系，在其中，不受外力的物体都保持静止或匀速直线运动。
\item[第二定律]施加于物体的合外力等于此物体的质量与加速度的乘积。
\item[第三定律]当两个物体互相作用时，彼此施加于对方的力，其大小相等、方向相反。
\end{description}
\end{block}
\end{frame}

\begin{frame}{建模方法}
\begin{block}{热力学：热力学定律}
\begin{description}
\item[第零定律]若两个热力学系统均与第三个系统处于热平衡状态，此两个系统也必互相处于热平衡。
\item [第一定律]物体内能的增加等于物体吸收的热量和对物体所作的功的总和.（系统经过绝热循环，其所做的功为零，因此第一类永动机是不可能的（即不消耗能量做功的机械）。
\item [第二定律]孤立系统自发地朝着热力学平衡方向──最大熵状态──演化，同样地，第二类永动机永不可能实现。（不可能从单一热源吸取热量，使之完全变为有用功而不产生其他影响）
\item [第三定律]热力学系统的熵在温度趋近于绝对零度时趋于定值，特别地，对于完整晶体，这个定值为零
\end{description}
\end{block}
\end{frame}

\begin{frame}{黑箱建模（实验法）}
\begin{description}
\item[系统辨识]给系统添加测试信号，记录输出响应，使用数学模型逼近。
\item[神经网络]训练网络权值，最小化输出的最小二乘误差。神经网络建模可看做系统辨识的特例。
\end{description}
系统辨识已经形成了一门独立的学科。本章重点研究分析法。
\end{frame}

\subsection{控制理论中的系统模型}
\begin{frame}{控制理论中的系统模型}
\begin{block}{数学模型}
系统内各变量之间的关系\\
静态$\longrightarrow$各阶导数为零\\
动态$\longrightarrow$各阶变量之间关系
\end{block}
\begin{block}{数学模型的类型}
\begin{description}
\item[时域($t$)]微分方程，差分方程，状态方程
\item[复数域($s$)]传递函数，结构图
\item[频率域($\omega$)]频率特性
\end{description}
\end{block}
\end{frame}

\section[控制系统的数学模型]{控制系统的数学模型}
\subsection{控制系统的微分方程}
\begin{frame}{控制系统的微分方程}
\begin{block}{线性定常系统微分方程的一般形式~~~~\tikz{\path[draw,-stealth'] (0,0)--node[above]{$r(t)$} (2em,0) coordinate(t);\node[anchor=west,draw,rectangle] (b) at (t) {System};
\path[draw,-stealth'] (b.east) -- node[above]{$c(t)$}++(2em,0);}}

\begin{align}
	&a_{n}\dod[n]{c(t)}{t}+a_{n-1}\dod[n-1]{c(t)}{t}+\cdots+a_{1}\dod{c(t)}{t}+a_0c(t)\nonumber\\
	=&b_{m}\dod[m]{r(t)}{t}+b_{m-1}\dod[m-1]{r(t)}{t}\cdots+b_{1}\dod{r(t)}{t}+b_0r(t)\label{eq:difEq}
\end{align}
\end{block}
\begin{block}{建立微分方程的一般步骤}
\begin{enumerate}
\item 确定输入量、输出量和扰动量，并根据需要引进一些中间变量。
\item 根据物理或化学定律，列出微分方程。
\item 消去中间变量后得到描述输出量与输入量(包括扰动量)关系的微分方程（标准形式）。
\end{enumerate}
\end{block}
\end{frame}
%\note{首先根据物理、化学等相关定律，列写输入变量与输出变量之间关系的微分方程。注意将输入与输出分别写在等式的两端。输入量在等号的右边，输出量在等号的左边，均按降次排列}

\subsection{建模实例}
\begin{frame}{\example{R-L-C 串连电路}{ex:RLC}}
\begin{columns}[T]
    \column{.5\textwidth}

由基尔霍夫定律得
\begin{align*}
&u_r(t)=L\dod{i(t)}{t}+Ri(t)+u_c(t)\\
&i(t)=C\dod{u_c(t)}{t}
\end{align*}
消去中间变量$i(t)$
\begin{align*}
&LC\dod[2]{u_c(t)}{t}+RC\dod{u_c(t)}{t}+u_c(t)=u_r(t)
\end{align*}

    \column[T]{.5\textwidth}
    \begin{center}
        \includegraphics{../figs/L3p1.pdf}
    \end{center}
\end{columns}
\end{frame}

\begin{frame}{\example{有负载效应的电路}{ex:load}}
%\begin{columns}[T]
%\column{.5\textwidth}
\vspace{-1.5em}
\begin{center}
\includegraphics{../figs/L3p2.pdf}
\end{center}
	\begin{align*}
	&u_r=\frac{1}{C_1}\int(i_1-i_2)\dif t+i_1R_1\\
	&\frac{1}{C_1}\int(i_1-i_2)\dif t=i_2R_2+u_c\\
	&u_c=\frac{1}{C_2}\int i_2\dif t
	\end{align*}
	消去中间变量$i_1$，$i_2$，整理得
	\begin{align*}
	&R_1R_2C_1C_2\dod[2]{u_c}{t}+(R_1C_1+R_2C_2+R_1C_2)\dod{u_c}{t}+u_c=u_r
	\end{align*}
\end{frame}
\note{负载效应，两个RC网络串联，其最终的传递函数不一定等于两个RC网络传递函数的乘积，因为后续网络会产生负载效应\\
推导过程：
$$u_c=\frac{1}{C_2}\int i_2dt\Rightarrow i_2=C_2u'_c$$
$$\frac{1}{C_1}\int(i_1-i_2)dt=i_2R_2+u_c=i_2R_2+u_c\Rightarrow i_1-i_2=C_1R_2i'_2+C_1u'_c$$
1代入2
$$i_1=C_1C_2R_2u''_c+C_1u'_c+C_2u'_c$$
\begin{align*}
u_r&=\frac{1}{C_1}\int(i_1-i_2)dt+i_1R_1\\
&=\frac{1}{C_1}\int[C_1C_2R_2u''_c+(C_1+C_2)u'_c-C_2u'_c]dt+i_1R_1\\
&=C_2R_2u'_c+u_c+C_1C_2R_1R_2u''_c+C_1R_1u'_c+C_2R_1u'_c\\
&=C_1C_2R_1R_2u''_c+(C_2R_2+C_1R_1+C_2R_1)u'_c+u_c
\end{align*}}

\begin{frame}{\example{机械位移系统}{ex:machine}}
\begin{columns}[t]
    \column{.5\textwidth}
由牛顿第二定律：
$f_M(t)=M\dod[2]{x(t)}{t}$\\
阻尼器模型：
$f_B(t)=B\dod{x(t)}{t}$\\
由胡克定律：
$f_K(t)=Kx(t)$\\
\begin{align*}
F(t)&=f_M(t)+f_B(t)+f_K(t)\\
&=M\dod[2]{x(t)}{t}+B\dod{x(t)}{t}+Kx(t)
\end{align*}
    \column[T]{.5\textwidth}
    \begin{center}
        \includegraphics{../figs/L3p3.pdf}
    \end{center}
\end{columns}
\end{frame}

\begin{frame}{\example{电枢控制式直流电动机}{ex:motor}}
\begin{columns}[t]
    \column{.65\textwidth}

\begin{eqnarray*}
\text{电枢回路}&u_r=Ri+E_b&\text{克希霍夫}\\
\text{电枢反电势}&E_b=c_e\cdot\omega_m&\text{楞次定律}\\
\text{电磁力矩}&M_m=c_mi&\text{安培定律}\\
\text{力矩平衡}&J_m\dot{\omega}_m+f_m\omega_m=M_m&\text{牛顿定律}\\
&\omega_m=\dot{\theta}&
\end{eqnarray*}
消去中间变量$i$，$M_m$，$E_b$可得：
$$\begin{matrix}T_m\dot{\omega}_m+\omega_m=K_mu_r\\
T_m\ddot{\theta}_m+\dot{\theta}_m=K_mu_r\end{matrix}~~
\begin{cases}T_m=J_mR/(R\cdot f_m+c_e\cdot c_m)&\text{电机时间常数 
}\\ K_m=c_m/(R\cdot f_m+c_e\cdot c_m)&\text{电机传递系数
}\end{cases}$$

    \column[T]{.35\textwidth}
    \begin{center}
        \includegraphics{../figs/L3p4.pdf}
    \end{center}
\end{columns}
\end{frame}

%\note{$E_b$为正比于电机转速的反电势，$c_e$为反电势常数，$c_m$为转矩常数，$J_m$为电动机的转动惯量，$f_m$为摩擦引起的阻尼系数直流电机，位置控制，无负载网络，位移系统。其中位置控制也可以叫做伺服系统，典型的伺服系统有救援机甲，手术机器人等。TED中的七个仿生机器人有一种蛇形机器人也是伺服系统。影视作品中有环太平洋，EVA，高达，苹果核战记等}

%:Lecture 2: 控制系统的复数域数学模型
\lecture{控制系统的复数域数学模型}{Lecture 2}
\section{控制系统的复数域数学模型}
\subsection{非线性模型的线性化}
\begin{frame}{非线性模型的线性化}
$$\ddot{c}+\dot{c}=\ddot{r}+r\longrightarrow\text{线性}$$
$$c\ddot{c}+\dot{c}=\ddot{r}+r\longrightarrow\text{非线性}$$
在给定工作点附近将非线性模型展开为泰勒级数
$$y=f(x)=f(x_0)+\left.\dod{f}{x}\right|_{x=x_0}(x-x_0)+\frac{1}{2}\left.\dod[2]{f}{x}\right|_{x=x_0}(x-x_0)^2+\cdots$$
略去2次及高阶项，得到线性化方程
$$y=y_0+K(x-x_0)$$
其中
$$y_0=f(x_0),~~~~K=\left.\dod{f}{x}\right|_{x=x_0}$$
\note{
关键看是否满足叠加原理\\
本节不作要求，但考研要求\\
推荐教材：斯坦福}
\end{frame}

\subsection{拉普拉斯变换}
\begin{frame}{导入}
	\begin{block}{关键词}
		法国\pause，18世纪后期\pause，（1782年《四库全书》，1787年台湾天地会反清，1796年白莲教起义）\pause，数学\pause，3L\pause（Lagrange，Laplace，Legendre）
	\end{block}
	\centering
	\begin{tabular}{ccc}
		\includegraphics[height=4cm]{../figs/C2Lagrange.jpg}&\includegraphics[height=4cm]{../figs/C2Laplace.jpg}&\includegraphics[height=4cm]{../figs/C2Legendre.jpg}\\
		拉格朗日&拉普拉斯&勒让德
	\end{tabular}
\end{frame}

\begin{frame}{导入}
	\centering
	\begin{tikzpicture}[x=1em,y=1em]
		\onslide<4->{\node[label=拉格朗日] (Lagrange) at (0,0) {\includegraphics[height=2cm]{../figs/C2Lagrange.jpg}};}
		\node[label=拉普拉斯] (Laplace) at (0,-9) {\includegraphics[height=2cm]{../figs/C2Laplace.jpg}};
		\onslide<3->{\node[label=泊松] (Poisson) at (8,-9) {\includegraphics[height=2cm]{../figs/C2Poisson.jpg}};}
		\onslide<5->{\node[label=傅里叶] (Fourier) at (8,0) {\includegraphics[height=2cm]{../figs/C2Fourier.jpg}};}
		\onslide<6->{\node[label=欧拉] (Euler) at (-8,0) {\includegraphics[height=2cm]{../figs/C2Euler.jpg}};}
		\onslide<2->{\node[label=达朗贝尔] (Alembert) at (-8,-9) {\includegraphics[height=2cm]{../figs/C2Alembert.jpg}};}
		\onslide<7->{\node[label=约翰·伯努利] (Johann) at (-16,0) {\includegraphics[height=2cm]{../figs/C2Johann_Bernoulli.jpg}};}
		\onslide<8->{\node[label=雅各布·伯努利] (Jakob) at (-16,-9) {\includegraphics[height=2cm]{../figs/C2Jakob_Bernoulli.jpg}};}
		\onslide<3->{\path[arr] (Laplace) -- (Poisson) node[midway,above] {博士生};}
		\onslide<4->{\path[arr] (Lagrange) -- (Poisson) node[midway,right] {博士生};}
		\onslide<5->{\path[arr] (Lagrange) -- (Fourier) node[midway,above] {博士生};}
		\onslide<6->{\path[arr] (Euler) -- (Lagrange) node[midway,above] {博士生};}
		\onslide<2->{\path[arr] (Alembert) -- (Laplace) node[midway,above] {博士生};}
		\onslide<7->{\path[arr] (Johann) -- (Euler) node[midway,above] {博士生};}
		\onslide<8->{\path[arr] (Jakob) -- (Johann) node[near end,right] {博士生};}
	\end{tikzpicture}
\end{frame}

\begin{frame}{第四讲学习目标}
	\begin{enumerate}
		\item 掌握拉普拉斯变换的定义
		\item  理解拉普拉斯变换微分定理的含义（对导数取拉普拉斯变换）
		\item 了解拉普拉斯变换积分定理
		\item 理解初值定理的含义
		\item 掌握终值定理的应用（求稳态）及其内涵
	\end{enumerate}
\end{frame}

\begin{frame}{拉普拉斯变换}
	\begin{block}{前测：配对定理}
		线性定理，微分定理，积分定理，初值定理，终值定理，位移定理，相似定理，卷积定理
	\end{block}
	\begin{block}{拉普拉斯变换的定义}
		$$F(s) = \int_0^\infty f(t)e^{-st}\dif t$$
	\end{block}
	其中$s$为复变量，上面的变换简称\emph{拉氏变换}，记作
	$$F(s) = \mathcal{L}[f(t)]$$
\end{frame}

\begin{frame}{拉普拉斯变换定理}
	\begin{block}{线性定理}
		$$\mathcal{L}[af_1(t)+bf_2(t)]=a\mL[f_1(t)]+b\mL[f_2(t)]=aF_1(s)+bF_2(s)$$
	\end{block}
	\onslide<2->{
	\begin{block}{微分定理\onslide<5->{——零初始条件}}
		\begin{align*}
			\mL\left[\fracd{f(t)}{t}\right]	&=sF(s)\onslide+<-4>{-f(0)}\\
			\onslide<3->{\mL\left[\fracd{^2f(t)}{t^2}\right]	&=s^2F(s)\onslide+<-4>{-[sf(0)+\dot{f}(0)]}\\
			&\cdots\\}
			\onslide<4->{\mL\left[\fracd{^nf(t)}{t^n}\right]	&=s^nF(s)\onslide+<-4>{-[s^{n-1}f(0)+s^{n-2}\dot{f}(0)+\cdots+f^{(n-1)}(0)]}}
		\end{align*}
	\end{block}}
\end{frame}

\begin{frame}{拉普拉斯变换定理}
	\begin{block}{积分定理\onslide<4->{——零初始条件}}
		\begin{align*}
			\mL\left[\int f(t)\dif t\right]	&=\frac{1}{s}F(s)\onslide+<-3>{+\frac{1}{s}f^{(-1)}(0)}\\
			\onslide<2->{\mL\left[\iint f(t)(\dif t)^2\right]	&=\frac{1}{s^2}F(s)\onslide+<-3>{+\frac{1}{s^2}f^{(-1)}(0)+\frac{1}{s}f^{(-2)}(0)}}\\
			\onslide<3->{&\cdots\\
			\mL\left[\idotsint f(t)(\dif t)^n\right]	&=\frac{1}{s^n}F(s)\onslide+<-3>{+\frac{1}{s^n}f^{(-1)}(0)+\cdots+\frac{1}{s}f^{(-n)}(0)}}
		\end{align*}
	\end{block}
	*$f^{(-1)}(0)=\left.\int f(t)\dif t\right|_{t=0}$
\end{frame}

\begin{frame}{拉普拉斯变换定理}
	\begin{block}{初值定理}
		$$f(0_+)=\limit{t}{0+}f(t)=\limit{s}{\infty}sF(s)$$
	\end{block}
	\pause
	\begin{block}{终值定理}
		$$f(\infty)=\limit{t}{\infty}f(t)=\limit{s}{0}sF(s)$$
	\end{block}
	\pause
	\begin{block}{位移定理}
		$$\mL[f(t-\tau_0)]=e^{-\tau_0s}F(s)$$
		$$\mL[e^{at}f(t)]=F(s-a)$$
	\end{block}
\end{frame}

\begin{frame}{拉普拉斯变换定理}
	\begin{block}{相似定理}
		$$\mL\left[f(\frac{t}{a})\right]=aF(as)$$
	\end{block}
	\pause
	\begin{block}{卷积定理}
		设$F_1(s)=\mL[f_1(t)]$，$F_2(s)=\mL[f_2(t)]$，则有
		$$F_1(s)F_2(s)=\mL\left[\int_0^tf_1(t-\tau)f_2(\tau)\dif\tau\right]$$
	\end{block}
\end{frame}

\begin{frame}{拉普拉斯反变换}
	$$f(t)=\mathcal{L}^{-1}[F(s)] = \frac{1}{2\pi}\int_{c-j\infty}^{c+j\infty}F(s)e^{st}\dif s,\quad (t>0)$$
	\begin{block}{\onslide<2->{如何求解下式的时间域函数？}\onslide<3->{——部分分式法}}
		\onslide<2->{$$F(s) = \frac{N(s)}{D(s)} = \frac{b_0s^m+b_1s^{m-1}+\cdots + b_{m-1}s+b_m}{s^n+a_1s^{n-1}+\cdots + a_{n-1}s+a_n}$$}
	\end{block}
	\onslide<4->{
	\begin{block}{首先，分母因式分解}
		$$F(s) = \frac{N(s)}{D(s)} = \frac{b_0s^m+b_1s^{m-1}+\cdots + b_{m-1}s+b_m}{(s-s_1)(s-s_2)\cdots(s-s_n)}$$
		其中$s_1$、$s_2$是$D(s)=0$的根。
	\end{block}}
\end{frame}

\begin{frame}{部分分式法——$D(s)=0$无重根}
	\begin{align*}
		F(s)	&=\frac{N(s)}{D(s)} = \frac{b_0s^m+b_1s^{m-1}+\cdots + b_{m-1}s+b_m}{(s-s_1)(s-s_2)\cdots(s-s_n)}\\
			&=\frac{c_1}{s-s_1}+\frac{c_2}{s-s_2}+\cdots+\frac{c_n}{s-s_n}=\sum_{i=1}^{n}\frac{c_i}{s-s_i}
	\end{align*}
	$c_1$待定，解法：
	\pause
	$$c_i=\limit{s}{s_i}(s-s_i)F(s)\quad\text{或}\quad c_i=\left.\frac{N(s)}{\dot{D}(s)}\right|_{s=s_i}$$
	*$\dot{D}(s)=\fracd{D(s)}{s}$
	$$f(t)=\mathcal{L}^{-1}[F(s)]=\mathcal{L}^{-1}\left[\sum_{i=1}^{n}\frac{c_i}{s-s_i}\right]=\sum_{i=1}^n c_ie^{s_it}$$
\end{frame}

\begin{frame}{\example{求$F(s)$的原函数}{ex:ReverseLaplace1}}
	$$F(s)=\frac{s^2+2s+2}{(s+1)(s+2)(s+3)}$$
	\begin{align*}
		\onslide<2->{F(s)	&=\frac{c_1}{(s+1)}+\frac{c_2}{(s+2)}+\frac{c_3}{(s+3)}}\\
		\onslide<3->{c_1	&=\limit{s}{-1}(s+1)F(s)=\limit{s}{-1}\frac{s^2+2s+2}{(s+2)(s+3)}=\frac{1}{2}}\\
		\onslide<4->{c_2	&=\limit{s}{-2}(s+2)F(s)=-2}\\
		\onslide<5->{c_3	&=\limit{s}{-3}(s+3)F(s)=\frac{5}{2}}\\
		\onslide<6>{f(t)=	&\mathcal{L}^{-1}[F(s)]=\mathcal{L}^{-1}\left[\sum_{i=1}^{n}\frac{c_i}{s-s_i}\right]=\frac{1}{2}e^{-t}-2e^{-2t}+\frac{5}{2}e^{-3t}\quad (t\ge 0)}
	\end{align*}
\end{frame}

\begin{frame}{\example{求$F(s)$的原函数}{ex:ReverseLaplace2}}
	$$F(s)=\frac{s-3}{s^2+2s+2}$$
	\begin{align*}
		F(s)	&=\frac{s-3}{(s+1)^2+1}=\frac{s+1}{(s+1)^2+1}-\frac{4}{(s+1)^2+1}\\
		f(t)	&=e^{-t}\cos t-4e^{-t}\sin t
	\end{align*}
\end{frame}

\begin{frame}{部分分式法——$D(s)=0$有重根}
	\begin{align*}
		F(s)	&=\frac{N(s)}{(s-s_1)^r(s-s_{r+1})\cdots(s-s_n)}\\
			&=\frac{c_r}{(s-s_1)^r}+\frac{c_{r-1}}{(s-s_1)^{r-1}}+\cdots\frac{c_1}{(s-s_1)}+\frac{c_{r+1}}{(s-s_{r+1})}+\cdots+\frac{c_n}{(s-s_n)}\\
		c_r	&=\limit{s}{s_1}(s-s_1)^rF(s)\\
		c_{r-1}&=\limit{s}{s_1}\dod{}{s}[(s-s_1)^rF(s)]\\
			&\cdots\\
		c_{r-j}&=\frac{1}{j!}\limit{s}{s_1}\dod[j]{}{s}[(s-s_1)^rF(s)]\\
			&\cdots\\
		c_{1}&=\frac{1}{(r-1)!}\limit{s}{s_1}\dod[r-1]{}{s}[(s-s_1)^rF(s)]
	\end{align*}
\end{frame}

\begin{frame}{部分分式法——$D(s)=0$有重根}
	有重根的部分分式法接拉普拉斯反变换可以总结为
	\begin{align*}
		f(t)	&=\mathcal{L}^{-1}[F(s)]\\
			&=\mathcal{L}^{-1}\left[\frac{c_r}{(s-s_1)^r}+\frac{c_{r-1}}{(s-s_1)^{r-1}}+\cdots+\frac{c_1}{s-s_1}+\frac{c_{r+1}}{s-s_{r+1}}+\cdots+\frac{c_n}{s-s_n}\right]\\
			&=\left[\frac{c_r}{(r-1)!}t^{r-1}+\frac{c_{r-1}}{(r-2)!}t^{r-2}+\cdots+c_2t+c_1\right]e^{s_1t}+\sum_{i=r+1}^n c_ie^{s_it}
	\end{align*}
\end{frame}

\begin{frame}{\example{求$F(s)$的原函数}{ex:ReverseLaplace3}}
	$$F(s)=\frac{s+2}{(s-1)^2(s+1)}$$
	\begin{align*}
		F(s)	&=\frac{c_1}{(s-1)^2}+\frac{c_2}{s-1}+\frac{c_3}{s+1}\\
		c_1	&=\limit{s}{1}(s-1)^2F(s)=\limit{s}{1}\frac{s+2}{s+1}=\frac{3}{2}\\
		c_2	&=\limit{s}{1}\fracd{}{s}[(s-1)^2F(s)]=-\frac{1}{4}\\
		c_3	&=\limit{s}{-1}(s+1)F(s)=\frac{1}{4}\\
		f(t)	&=e^t\left(\frac{3}{2}t-\frac{1}{4}\right)+\frac{1}{4}e^{-t}\quad (t\ge 0)
		\end{align*}
\end{frame}

\subsection{控制系统的复数域数学模型}
\begin{frame}{控制系统的复数域数学模型}
传递函数：经典控制的基础~~~~~~~~
\tikz{\draw[-stealth'](0,0)--node[above]{$r(t)$}(2em,0);
\node[draw,rectangle,anchor=west] (b) at (2em,0){$\ddot{c}+\dot{c}=\ddot{r}+r$};
\draw[-stealth'](b.east)--node[above]{$c(t)$}++(2em,0);}
\begin{block}{定义}
零初始条件下输出的拉氏变换比输入拉氏变换。
\end{block}
\begin{block}{典型外作用信号}
\begin{description}
\item[阶跃]$f(t)=R\cdot1(t)$
\item[斜坡]$f(t)=Rt$
\item[脉冲]$f(t)=\limit{t_0}{0}\frac{A}{t_0}[1(t)-1(t-t_0)]$
\end{description}
\end{block}
复数域下，理想脉冲信号表示为$\delta(t)$，且有：
$$\mathcal{L}[\delta(t)]=1$$
\end{frame}

\begin{frame}{微分方程与传递函数}
回忆线性定常微分方程的一般形式\eqref{eq:difEq}：
\begin{align*}
	&a_{0}\dod[n]{c(t)}{t}+a_{1}\dod[n-1]{c(t)}{t}+\cdots+a_{n-1}\dod{c(t)}{t}+a_nc(t)\nonumber\\
	=&b_{0}\dod[m]{r(t)}{t}+b_{1}\dod[m-1]{r(t)}{t}\cdots+b_{m-1}\dod{r(t)}{t}+b_mr(t)\end{align*}

零初始条件下，对上式两边取拉氏变换，得：
\begin{align}
(a_0s^n+a_1s^{n-1}+\cdots+a_{n-1}s+a_n)C(s)\nonumber\\
=(b_0s^m+b_1s^{m-1}+\cdots+b_{m-1}s+b_m)R(s)
\end{align}
该系统的传递函数定义为输出与输入之比：
\begin{equation}
G(s)=\frac{C(s)}{R(s)}=\frac{b_0s^m+b_1s^{m-1}+\cdots+b_{m-1}s+b_m}{a_0s^n+a_1s^{n-1}+\cdots+a_{n-1}s+a_n}
\end{equation}
\end{frame}

\begin{frame}{线性系统的数学模型}
	\centering
	\begin{tikzpicture}[x=1em,y=1em]
		\node[block] (sys) at (0,0) {线性系统};
		\node[block] (tf) at (-12,-8) {传递函数};
		\node[block] (de) at (0,-8) {微分方程};
		\node[block] (fc) at (12,-8) {频率特性};
		\path[arr] (sys) -- (tf);
		\path[arr] (sys) -- (de);
		\path[arr] (sys) -- (fc);
		\path[arr] (de) -- (tf) node[midway,above] {拉氏变换} node[midway,below] {\includegraphics[height=3cm]{../figs/C2Laplace.jpg}};
		\path[arr] (de) -- (fc) node[midway,above] {傅里叶变换} node[midway,below] {\includegraphics[height=3cm]{../figs/C2Fourier.jpg}};
	\end{tikzpicture}
\end{frame}

\begin{frame}{传递函数的零极点形式}
\begin{align}
G(s)&=\frac{C(s)}{R(s)}=\frac{b_0s^m+b_1s^{m-1}+\cdots+b_{m-1}s+b_m}{a_0s^n+a_1s^{n-1}+\cdots+a_{n-1}s+a_n}\nonumber\\
&=\frac{N(s)}{D(s)}=\frac{b_0(s-z_1)(s-z_2)\cdots(s-z_m)}{a_0(s-p_1)(s-p_2)\cdots(s-p_n)}
\end{align}
$s=z_i(i=1,2,\cdots,m)$是$N(s)=0$的根，称为传递函数的零点；$s=p_i(i=1,2,\cdots,n)$是$D(s)=0$的根，称为传递函数的极点。
\end{frame}

\begin{frame}{传递函数性质}
\begin{itemize}
\item TF与信号无关，只取决于系统的结构和参数，与系统的输入和输出位置有关
\item 仅适用于线性定常系统
\item 分母阶次$n$大于分子阶次$m$，$n>m$
\item 分母中$s$的最高阶次$n$即为系统的阶次
\item 不能反映非零初始条件的运动
\item 由相应零极点组成
\item 只反映单输入与单输出的 关系，不反映内部（多入多出系统使用TF矩阵）
\item 特征多项式$D(s)=a_0s^n+a_1s^{n-1}+\cdots+a_{n-1}s+a_n$
\item 特征方程$D(s)=0$
\item $D(s)=0$的根称为特征根。
\end{itemize}
\end{frame}

\subsection{微分方程求传递函数举例}
\begin{frame}{微分方程求传递函数举例}
由\exref{ex:RLC}，所得系统微分方程为：
$$LC\od[2]{u_c(t)}{t}+RC\od{u_c(t)}{t}+u_c(t)=u_r(t)$$
拉氏变换得：$LCs^2U_c(s)+RCsU_c(s)+U_c(s)=U_r(s)$，整理得
$$\frac{U_c(s)}{U_r(s)}=\frac{1}{LCs^2+RCs+1}$$
由\exref{ex:load}，所得系统微分方程为
$$R_1R_2C_1C_2\od[2]{u_c}{t}+(R_1C_1+R_2C_2+R_1C_2)\dod{u_c}{t}+u_c=u_r$$
拉氏变换得并整理得：
$$\frac{U_c(s)}{U_r(s)}=\frac{1}{R_1R_2C_1C_2s^2+(R_1C_1+R_2C_2+R_1C_2)s+1}$$
\end{frame}

\begin{frame}{微分方程求传递函数举例}
由\exref{ex:machine}，所得系统微分方程为：
$$M\od[2]{x(t)}{t}+B\od{x(t)}{t}+Kx(t)=F(t)$$
拉氏变换并整理得：
$$\frac{X(s)}{F(s)}=\frac{1}{Ms^2+Bs+K}$$
由\exref{ex:motor}，所得系统微分方程为
$$T_m\dot{\omega}_m+\omega_m=K_mu_r\text{~~或~~}T_m\ddot{\theta}_m+\dot{\theta}_m=K_mu_r$$
拉氏变换得并整理得：
$$\frac{\Omega_m(s)}{U_r(s)}=\frac{K_m}{T_ms+1}\text{~~或~~}\frac{\Theta_m(s)}{U_r(s)}=\frac{K_m}{s(T_ms+1)}$$
\end{frame}

\subsection{典型环节的数学模型}
\begin{frame}{典型环节的数学模型}
\begin{block}{1）比例环节}
其输出量和输入量的关系，
$$y(t)=Kr(t)$$
式中$K$为环节的放大系数，为常数。传递函数为
$$G(s)=\frac{Y(s)}{R(s)}=K$$
\begin{description}
\item[特点：]输出与输入量成比例，无失真和时间延迟。
\item[实例：]电子放大器，齿轮，电阻(电位器) 等。
\end{description}
\end{block}
\end{frame}

\begin{frame}{典型环节的数学模型}
\begin{block}{2）惯性环节}
微分方程式表示
$$T\dod{y(t)}{t}+y(t)=r(t)$$
式中$T$为环节的时间常数。传递函数为
$$G(s)=\frac{Y(s)}{R(s)}=\frac{1}{Ts+1}$$
\begin{description}
\item[特点：]对突变的输入,其输出不能立即复现，输出无振荡。
\item[实例：]RC网络，直流伺服电动机的传递函数也包含这一环节。
\end{description}
\end{block}
\end{frame}

\begin{frame}{典型环节的数学模型}
\begin{block}{3）积分环节}
微分方程式表示
$$y(t)=\int r(t)\dif t$$
传递函数为
$$G(s)=\frac{Y(s)}{R(s)}=\frac{1}{s}$$
\begin{description}
\item[特点：]输出量与输入量的积分成正比例，当输入消失，输出具有记忆功能。
\item[实例：]模拟计算机中的积分器。
\end{description}
\end{block}
\end{frame}

\begin{frame}{典型环节的数学模型}
\begin{block}{4）微分环节}
是积分的逆运算，微分方程式表示
$$y(t)=\tau\od{r(t)}{t}$$
式中$\tau$为环节的时间常数，传递函数为
$$G(s)=\frac{Y(s)}{R(s)}=\tau s$$
\begin{description}
\item[特点：]输出量正比输入量变化的速度，能预示输入信号的变化趋势。
\item[实例：]测速发电机输出电压与输入角度间的传递函数即为微分环节。
\end{description}
\end{block}
\end{frame}

\begin{frame}{典型环节的数学模型}
\begin{block}{5）振荡环节：}
二阶微分方程式来表示
$$T^2\od[2]{y(t)}{t}+2\zeta T\od{y(t)}{t}+y(t)=r(t)$$
传递函数为
$$G(s)=\frac{Y(s)}{R(s)}=\frac{1}{T^2s^2+2\zeta Ts+1}$$
\begin{description}
\item[特点：]环节中有两个独立的储能元件，并可进行能量交换。
\item[实例：]RLC电路的输出与输入电压间的传递函数。
\end{description}
\end{block}
\end{frame}

\subsection{典型环节的特点}
\begin{frame}{典型环节的特点}
\begin{itemize}
\item 环节：具有相同形式传递函数的元部件的分类。
\item 不同的元部件可以有相同的传递函数；
\item 若输入输出变量选择不同，同一部件可以有不同的传递函数 ；
\item 任一传递函数都可看作典型环节的组合，即传递函数的典型环节形式。
\end{itemize}
例如：
$$G(s)=\frac{K(2s+1)}{s^v(Ts+1)(\tau^2s^2+2\epsilon\tau s+1)}$$
\end{frame}

\begin{frame}{结构相似系统}
许多表面上看来似乎毫无共同之处的控制系统，其运动规律可能完全一样，可以用一个运动方程来表示，称它们为结构相似系统。
  
\exref{ex:RLC}的RLC电路和\exref{ex:machine}机械平移系统就可以用同一个数学表达式分析，具有相同的数学模型。
$$LC\od[2]{u_c(t)}{t}+RC\od{u_c(t)}{t}+u_c(t)=u_r(t)$$
$$M\od[2]{x(t)}{t}+B\od{x(t)}{t}+Kx(t)=F(t)$$
\end{frame}

%:Lecture 3: 动态结构图及等效变换
\lecture{动态结构图及等效变换}{Lecture 3}
\section{动态结构图及等效变换}
\subsection{动态结构图的组成}
\begin{frame}{动态结构图的组成}
\begin{block}{1、信号线：有箭头的直线，箭头表示信号传递方向。}
\begin{center}
\tikz{\draw[-stealth',thick](0,0)--node[above,midway]{$U(s)$}(8em,0);}
\end{center}
\end{block}
\vspace{2em}
\begin{block}{2、引出点：信号引出或测量的位置。}
\begin{columns}[t]
\column{.4\textwidth}\tikz{\draw[-stealth',thick](0,0)--node[above,at start]{$U(s)$} node[above,very near end]{$U(s)$}(8em,0);\draw[-stealth',thick](3em,0)--node[right,near end]{$U(s)$}(3em,-4em);}
\column{.4\textwidth}从同一信号线上引出的信号，数值和性质完全相同
\end{columns}
\end{block}
\end{frame}

\begin{frame}{动态结构图的组成}
\begin{block}{3、综合点：对两个或两个以上的信号进行代数运算，“$+$”表示相加，常省略，“$-$”表示相减。}
\begin{center}
\tikz{
\draw[-stealth',thick](0,0)--node[above,very near start]{$R(s)$}(8em,0);
\draw (8.5em,0) circle (0.5em);
\draw (8.15em,0.35em)--(8.85em,-0.35em);
\draw (8.15em,-0.35em)--(8.85em,0.35em);
\draw[-stealth',thick](9em,0)--node[above,very near end]{$U(s)$}(15em,0);
\draw[-stealth',thick](8.5em,-3em)--node[right,very near start]{$B(s)$} node[left,very near end]{$\pm$}(8.5em,-0.5em);
}
\end{center}
\end{block}

\begin{block}{4、方框：表示典型环节或其组合，框内为对应的传递函数 ，两侧为输入、输出信号线。}
\begin{center}
\tikz{
\draw[-stealth',thick](0,0)--node[above,at start]{$R(s)$} ++(4em,0) coordinate(t);
\node[draw,rectangle,minimum width=4em,minimum height=3em,anchor=west] (b) at(t) {$G(s)$};
\draw[-stealth',thick](b.east)--node[above,near end]{$Y(s)=R(s)G(s)$}++(7em,0);}
\end{center}
\end{block}
\end{frame}

\begin{frame}{动态结构图的建立}
\begin{center}
\tikz{
\node[rectangle,draw] (b1)at(0,0) {工作原理图};
\draw[triangle 90 cap reversed-triangle 90 cap,line width=1ex,color=cyan](b1.north)|- ++(4em,4em)coordinate(t);
\node[rectangle,draw,anchor=west] (b2)at(t) {方框图};
\draw[triangle 90 cap reversed-triangle 90 cap,line width=1ex,color=cyan](b2.east)-| node[above]{$G_i(s)$}++(4em,-4em)coordinate(t);
\node[rectangle,draw,anchor=north] (b3)at(t) {系统结构图};
\draw[triangle 90 cap reversed-triangle 90 cap,line width=1ex,color=blue](b1.south)|- ++(4em,-4em)coordinate(t);
\node[rectangle,draw,anchor=west] (b4)at(t) {微分方程组};
\draw[triangle 90 cap reversed-triangle 90 cap,line width=1ex,color=blue](b4.east)-| node[below]{$L$}(b3.south);}
\end{center}
\end{frame}

\begin{frame}{\example{有负载效应的RC网络}{ex:loadRC}}
\begin{block}{1）由微分方程得到规范的因果关系}
\begin{equation*}
\begin{matrix}
\color{red}i_1(t)=\frac{u_r(t)-u_1(t)}{R_1}&\Longrightarrow&\color{red}[U_r(s)-U_{c_1}(s)]\frac{1}{R_1}=I_1(s)\\
\color{blue}u_1(t)=\frac{1}{C_1}\int(i_1(t)-i_2(t))dt&\Longrightarrow&\color{blue}[I_1(s)-I_2(s)]\frac{1}{sC_1}=U_1(s)\\
\color{green}i_2(t)=\frac{u_1(t)-u_C(t)}{R_2}&\Longrightarrow&\color{green}[U_1(s)-U_C(s)]\frac{1}{R_2}=I_2(s)\\
\color{cyan}u_C(t)=\frac{1}{C_2}\int i_2(t)dt&\Longrightarrow&\color{cyan}I_2(s)\frac{1}{sC_2}=U_C(s)\\
\end{matrix}
\end{equation*}
\end{block}
\begin{block}{2）绘动态结构图。按照变量的传递顺序，依次将各元件的结构图连接起来}
\begin{center}
\tikz{
\draw[-latex,thick](0,0)--node[above,very near start]{$U_r(s)$}(2em,0);
\draw[cross] (2.5em,0) circle (0.5em) ++(0,-0.5em)coordinate(c1);
\draw[-latex,thick](3em,0)--++(1.5em,0)coordinate(t);
\node[draw,rectangle,anchor=west] (b1) at (t) {$\frac{1}{R_1}$};
\draw[-latex,thick](b1.east)--node[above]{$I_1(s)$}++(3em,0)coordinate(t);
\draw[cross] (t)[xshift=0.5em] circle (0.5em)++(0.5em,0)coordinate(t)++(-0.5em,0.5em)coordinate(c2);
\draw[-latex,thick](t)--++(1.5em,0)coordinate(t);
\node[draw,rectangle,anchor=west] (b2) at (t) {$\frac{1}{sC_1}$};
\draw[-latex,thick](b2.east)--node[above,at end]{$U_{c_1}(s)$}++(1.5em,0)coordinate(o1)--++(1.5em,0)coordinate(t);
\draw[cross] (t)[xshift=0.5em] circle (0.5em)++(0.5em,0)coordinate(t)++(-0.5em,-0.5em)coordinate(c3);
\draw[-latex,thick](t)--++(1.5em,0)coordinate(t);
\node[draw,rectangle,anchor=west] (b3) at (t) {$\frac{1}{R_2}$};
\draw[-latex,thick](b3.east)--++(1em,0)coordinate(o2)--++(1em,0)coordinate(t);
\node[draw,rectangle,anchor=west] (b4) at (t) {$\frac{1}{sC_2}$};
\draw[-latex,thick](b4.east)--node[above,very near end]{$U_c(s)$}++(1.5em,0)coordinate(o3)--++(1.5em,0)coordinate(t);
\draw[-latex,thick](o1)--++(0,-3em)-|node[right,near end]{$U_{c_1}(s)$}node[left,very near end]{$-$}(c1);
\draw[-latex,thick](o2)--++(0,3em)-|node[right,very near end]{$-$}(c2);
\draw[-latex,thick](o3)--++(0,-3em)-|node[right,near end]{$U_c(s)$}node[left,very near end]{$-$}(c3);
\draw[dashed,color=red] (-1em,3em)rectangle(6.7em,-3.5em);
\draw[dashed,color=blue] (8em,3.5em)rectangle(11em,-2em);
\draw[dashed,color=green] (17em,2.5em)rectangle(22em,-3.5em);
\draw[dashed,color=cyan] (23em,2.5em)rectangle(26em,-2em);
}
\end{center}
\end{block}
\end{frame}

\subsection{典型连接方式及等效变换}
\begin{frame}{典型连接方式及等效变换}
\begin{block}{1、串联及等效}
\begin{center}
\tikz{
\draw[-latex,thick](0,0)--node[above,at start]{$X(s)$} ++(2em,0) coordinate(t);
\node[draw,rectangle,minimum width=4em,minimum height=3em,anchor=west] (b1) at(t) {$G_1(s)$};
\draw[-latex,thick](b1.east)--node[above,midway]{$X_1(s)$}++(3em,0)coordinate(t);
\node[draw,rectangle,minimum width=4em,minimum height=3em,anchor=west] (b2) at(t) {$G_2(s)$};
\draw[-latex,thick](b2.east)--node[above,midway]{$Y(s)$}++(3em,0)coordinate(t);
\draw[-latex,line width=1ex](8em,-2em)|-++(3em,-3em)coordinate(t);
\begin{scope}[xshift=13em,yshift=-5em]
\draw[-latex,thick](0,0)--node[above,at start]{$X(s)$} ++(2em,0) coordinate(t);
\node[draw,rectangle,minimum width=4em,minimum height=3em,anchor=west] (b1) at(t) {$G(s)$};
\draw[-latex,thick](b1.east)--node[above,midway]{$Y(s)$}++(3em,0)coordinate(t);
\end{scope}}
\end{center}
$$\frac{X_1(s)}{X(s)}=G_1(s),~~~\frac{Y(s)}{X_1(s)}=G_2(s)$$
$$G(s)=\frac{Y(s)}{X(s)}=G_1(s)G_2(s)$$
\end{block}
\end{frame}

\begin{frame}{典型连接方式及等效变换}
\begin{block}{2、并联及等效}
\begin{center}
\tikz{
\draw[-latex,thick](0,0)--node[above,at start]{$X(s)$} ++(2em,0) coordinate(o)--++(1em,0)coordinate(t);
\node[draw,rectangle,minimum width=4em,minimum height=2em,anchor=west] (b1) at(t) {$G_1(s)$};
\draw[-latex,thick](b1.east)-|node[above,midway]{$Y_1(s)$}++(3em,-1em)coordinate(t);
\draw[cross] (t)[yshift=-0.5em] circle (0.5em)++(0.5em,0)coordinate(t)++(-0.5em,-0.5em)coordinate(c);
\draw[-latex,thick](t)--node[above,near end]{$Y(s)$}++(3em,0)coordinate(t);
\draw[-latex,thick](o)|-++(1em,-3em)coordinate(t);
\node[draw,rectangle,minimum width=4em,minimum height=2em,anchor=west] (b2) at(t) {$G_2(s)$};
\draw[-latex,thick](b2.east)-|node[below]{$Y_2(s)$}node[right,very near end]{$\pm$}(c);
\draw[-latex,line width=1ex](15em,-2em)--++(2em,0);
\begin{scope}[xshift=19em,yshift=-2em]
\draw[-latex,thick](0,0)--node[above,at start]{$X(s)$} ++(2em,0) coordinate(t);
\node[draw,rectangle,minimum width=4em,minimum height=2em,anchor=west] (b3) at(t) {$G(s)$};
\draw[-latex,thick](b3.east)--node[above,midway]{$Y(s)$}++(3em,0)coordinate(t);
\end{scope}}
\end{center}
\begin{align*}
Y(s)&=Y_1(s)\pm Y_2(s)=X(s)G_1(s)\pm X(s)G_2(s)\\
&=X(s)[G_1(s)\pm G_2(s)]=X(s)G(s)
\end{align*}
$$G(s)=G_1(s)\pm G_2(s)$$
\end{block}
\end{frame}

\begin{frame}{典型连接方式及等效变换}
\begin{block}{3、反馈及等效}
\begin{center}
\tikz{
\draw[-latex,thick](0,0)--node[above,at start]{$R(s)$} ++(2em,0) coordinate(t);
\draw[cross] (t)[xshift=0.5em] circle (0.5em)++(0.5em,0)coordinate(t)++(-0.5em,-0.5em)coordinate(c);
\draw[-latex,thick](t)--node[above]{$E(s)$}++(3em,0)coordinate(t);
\node[draw,rectangle,minimum width=4em,minimum height=2em,anchor=west] (b1) at(t) {$G(s)$};
\draw[-latex,thick](b1.east)--++(2em,0)coordinate(o)--node[above,near end]{$Y(s)$}++(1em,0);
\draw[-latex,thick](o)|-++(-2em,-3em)coordinate(t);
\node[draw,rectangle,minimum width=4em,minimum height=2em,anchor=east] (b2) at(t) {$H(s)$};
\draw[-latex,thick](b2.west)-|node[right,near end]{$B(s)$}node[left,very near end]{$\mp$}(c);
\draw[-latex,line width=1ex](15em,-2em)--++(2em,0);
\begin{scope}[xshift=19em,yshift=-2em]
\draw[-latex,thick](0,0)--node[above,at start]{$R(s)$} ++(2em,0) coordinate(t);
\node[draw,rectangle,minimum width=4em,minimum height=2em,anchor=west] (b3) at(t) {$\frac{G(s)}{1\pm G(s)H(s)}$};
\draw[-latex,thick](b3.east)--node[above,midway]{$Y(s)$}++(3em,0)coordinate(t);
\end{scope}}
\end{center}
\begin{align*}
&Y(s)=E(s)G(s),~~E(s)=R(s)\mp B(s),~~B(s)=Y(s)H(s)\\
&Y(s)=[R(s)\mp B(s)]G(s)=R(s)G(s)\mp Y(s)H(s)G(s)\\
&Y(s)[1\pm H(s)G(s)]=R(s)G(s)\\
&\frac{Y(s)}{R(s)}=\frac{G(s)}{1\pm H(s)G(s)}
\end{align*}
\end{block}
\end{frame}

\subsection{等效移动规则}
\begin{frame}{等效移动规则}
\begin{block}{1、引出点的移动}
1）前移：在移动支路中串入所越过的传递函数方框
\begin{center}
\tikz{
\draw[-latex,thick](0,0)--node[below,at start]{$X_1$} ++(2em,0) coordinate(t);
\node[draw,rectangle,anchor=west] (b1) at(t) {$G(s)$};
\draw[-latex,thick](b1.east)--++(1em,0)coordinate(o)--node[above,near end]{$X_2$}++(1em,0);
\draw[-latex,thick](o)|-node[below,near end]{$X_2$}++(2em,-2em)coordinate(t);
\draw[-latex,line width=1ex](10em,-1em)--++(2em,0);
\begin{scope}[xshift=13em]
\draw[-latex,thick](0,0)--node[above,at start]{$X_1$} ++(1em,0)coordinate(o)--++(2em,0) coordinate(t);
\node[draw,rectangle,anchor=west] (b2) at(t) {$G(s)$};
\draw[-latex,thick](b2.east)--node[above,very near end]{$X_2$}++(3em,0);
\draw[-latex,thick](o)|-++(2em,-2em)coordinate(t);
\node[draw,rectangle,anchor=west] (b3) at(t) {$G(s)$};
\draw[-latex,thick](b3.east)--node[above,very near end]{$X_2$}++(3em,0);
\end{scope}}
\end{center}
2）后移：在移动支路中串入所越过的传递函数的倒数方框 
\begin{center}
\tikz{
\draw[-latex,thick](0,0)--node[above,at start]{$X_1$} ++(1em,0)coordinate(o)--++(2em,0) coordinate(t);
\node[draw,rectangle,anchor=west] (b1) at(t) {$G(s)$};
\draw[-latex,thick](b1.east)--node[above,very near end]{$X_2$}++(3em,0);
\draw[-latex,thick](o)|-node[above,at end]{$X_1$}++(8em,-2em)coordinate(t);
\draw[-latex,line width=1ex](11em,-1em)--++(2em,0);
\begin{scope}[xshift=14em]
\draw[-latex,thick](0,0)--node[above,at start]{$X_1$} ++(2em,0) coordinate(t);
\node[draw,rectangle,anchor=west] (b2) at(t) {$G(s)$};
\draw[-latex,thick](b2.east)--++(1em,0)coordinate(o)--node[above,very near end]{$X_2$}++(5em,0);
\draw[-latex,thick](o)|-++(2em,-2em)coordinate(t);
\node[draw,rectangle,anchor=west] (b3) at(t) {$1/G(s)$};
\draw[-latex,thick](b3.east)--node[above,very near end]{$X_1$}++(1em,0);
\end{scope}}
\end{center}
\end{block}
\end{frame}

\begin{frame}{等效移动规则}
\begin{block}{2、综合点的移动}
1）前移：在移动支路中串入所越过的传递函数的倒数方框 
\begin{center}
\tikz{
\draw[-latex,thick](0,0)--node[above,at start]{$X_1$} ++(2em,0) coordinate(t);
\node[draw,rectangle,anchor=west] (b1) at(t) {$G(s)$};
\draw[-latex,thick](b1.east)--++(1em,0)coordinate(t);
\draw[cross] (t)[xshift=0.5em] circle (0.5em)++(0.5em,0)coordinate(t)++(-0.5em,-0.5em)coordinate(c);
\draw[-latex,thick](t)--node[above,at end]{$X_2$}++(2em,0);
\draw[latex-,thick](c)--node[right,at start]{$-$}node[right,at end]{$X_3$}++(0,-2em);
\draw[-latex,line width=1ex](10em,-1em)--++(2em,0);
\begin{scope}[xshift=13em]
\draw[-latex,thick](0,0)--node[above,at start]{$X_1$} ++(1em,0) coordinate(t);
\draw[cross] (t)[xshift=0.5em] circle (0.5em)++(0.5em,0)coordinate(t)++(-0.5em,-0.5em)coordinate(c);
\draw[-latex,thick](t)--++(1em,0)coordinate(t);
\node[draw,rectangle,anchor=west] (b2) at(t) {$G(s)$};
\draw[-latex,thick](b2.east)--node[above,very near end]{$X_2$}++(3em,0);
\draw[latex-,thick](c)|-node[right,at start]{$-$}++(2em,-2em)coordinate(t);
\node[draw,rectangle,anchor=west] (b3) at(t) {$1/G(s)$};
\draw[latex-,thick](b3.east)--node[above,very near end]{$X_3$}++(3em,0);
\end{scope}}
\end{center}
2）后移：在移动支路中串入所越过的传递函数方框 
\begin{center}
\tikz{
\draw[-latex,thick](0,0)--node[above,at start]{$X_1$} ++(2em,0) coordinate(t);
\draw[cross] (t)[xshift=0.5em] circle (0.5em)++(0.5em,0)coordinate(t)++(-0.5em,-0.5em)coordinate(c);
\draw[-latex,thick](t)--++(2em,0)coordinate(t);
\draw[latex-,thick](c)--node[right,at start]{$-$}node[right,at end]{$X_3$}++(0,-2em);
\node[draw,rectangle,anchor=west] (b1) at(t) {$G(s)$};
\draw[-latex,thick](b1.east)--node[above,at end]{$X_2$}++(1em,0)coordinate(t);
\draw[-latex,line width=1ex](10em,-1em)--++(2em,0);
\begin{scope}[xshift=15em]
\draw[-latex,thick](0,0)--node[above,at start]{$X_1$} ++(2em,0) coordinate(t);
\node[draw,rectangle,anchor=west] (b2) at(t) {$G(s)$};
\draw[-latex,thick](b2.east)--++(1em,0)coordinate(t);
\draw[cross] (t)[xshift=0.5em] circle (0.5em)++(0.5em,0)coordinate(t)++(-0.5em,-0.5em)coordinate(c);
\draw[-latex,thick](t)--node[above,at end]{$X_2$}++(1em,0)coordinate(t);
\draw[latex-,thick](c)|-++(-2em,-2em)coordinate(t);
\node[draw,rectangle,anchor=east] (b3) at(t) {$G(s)$};
\draw[latex-,thick](b3.west)--node[above,very near end]{$X_3$}++(-2em,0);
\end{scope}}
\end{center}
\end{block}
\end{frame}

\begin{frame}{等效移动规则}
\begin{block}{3）相邻综合点移动}
\begin{center}
\tikz{
\draw[-latex,thick](0,0)--node[above,at start]{$X_1$} ++(2em,0) coordinate(t);
\draw[cross] (t)[xshift=0.5em] circle (0.5em)++(0.5em,0)coordinate(t)++(-0.5em,-0.5em)coordinate(c);
\draw[-latex,thick](t)--++(2em,0)coordinate(t);
\draw[latex-,thick](c)--node[right,at end]{$X_2$}++(0,-2em);
\draw[cross] (t)[xshift=0.5em] circle (0.5em)++(0.5em,0)coordinate(t)++(-0.5em,0.5em)coordinate(c);
\draw[-latex,thick](t)--node[above,at end]{$Y$}++(2em,0);
\draw[latex-,thick](c)--node[right,at end]{$X_3$}++(0,2em);
\draw[-latex,line width=1ex](10em,-1em)--++(2em,0);
\begin{scope}[xshift=13em]
\draw[-latex,thick](0,0)--node[above,at start]{$X_1$} ++(2em,0) coordinate(t);
\draw[cross] (t)[xshift=0.5em] circle (0.5em)++(0.5em,0)coordinate(t)++(-0.5em,0.5em)coordinate(c);
\draw[-latex,thick](t)--++(2em,0)coordinate(t);
\draw[latex-,thick](c)--node[right,at end]{$X_3$}++(0,2em);
\draw[cross] (t)[xshift=0.5em] circle (0.5em)++(0.5em,0)coordinate(t)++(-0.5em,-0.5em)coordinate(c);
\draw[-latex,thick](t)--node[above,at end]{$Y$}++(2em,0);
\draw[latex-,thick](c)--node[right,at end]{$X_2$}++(0,-2em);
\end{scope}}
\end{center}
\vspace{-1em}
\begin{itemize}
\item 相邻综合点之间可以随意调换位置
\item 相邻引出点之间可以随意调换位置 
\end{itemize}
\vspace{-2em}
\begin{center}
\tikz{
\node[draw=orange,rectangle callout,line width=2pt, callout relative pointer={(1.5,1)},minimum height=3em,minimum width=22em]{注意：相邻引出点和综合点之间不能互换!};}
\end{center}
\end{block}
\end{frame}

\subsection{结构图等效变换方法}
\begin{frame}{结构图等效变换方法}
\begin{enumerate}
\item 若有三种典型结构，串联，并联或反馈，则直接用公式先化简;
\item 若没有三种典型结构，且回路之间有交叉，则必须移位。
\item 比较点只能向相邻比较点方向移位，引出点只能向相邻引出点方向移位。且移位后，常常伴随着与相邻综合点或相邻引出点交换位置或综合，以解交叉。
\item 由内回路向外回路一层层简化。
\item \alert{注意：相邻引出点和综合点之间不能互换!}
\end{enumerate}
\end{frame}

\begin{frame}{\example{结构图简化：电枢控制直流电动机}{}}
\begin{center}
\tikz{
\draw[-latex,thick](0,0)--node[above,at start]{$U_a$} ++(2em,0) coordinate(t);
\draw[cross] (t)[xshift=0.5em] circle (0.5em)++(0.5em,0)coordinate(t)++(-0.5em,-0.5em)coordinate(c);
\draw[-latex,thick](t)--++(1em,0)coordinate(t);
\node[draw,rectangle,anchor=west] (b1) at(t) {$\frac{1}{R}$};
\draw[-latex,thick](b1.east)--++(1em,0)coordinate(t);
\node[draw,rectangle,anchor=west] (b2) at(t) {$C_m$};
\draw[-latex,thick](b2.east)--++(1em,0)coordinate(t);
\node[draw,rectangle,anchor=west] (b3) at(t) {$\frac{1}{J_ms+f_m}$};
\draw[-latex,thick](b3.east)--node[above,at end]{$\omega$}++(1em,0)coordinate(o)--++(1em,0)coordinate(t);
\node[draw,rectangle,anchor=west] (b4) at(t) {$\frac{1}{s}$};
\draw[-latex,thick](b4.east)--node[above,at end]{$\theta$}++(2em,0);
\draw[-latex,thick](o)|-++(-3em,-2em)coordinate(t);
\node[draw,rectangle,anchor=east] (b5) at(t) {$C_e$};
\draw[-latex,thick](b5.west)-|node[left,at end]{$-$}(c);
}
\vspace{1em}
\tikz{
\draw[-latex,thick](0,0)--node[above,at start]{$U_a$} ++(2em,0) coordinate(t);
\draw[cross] (t)[xshift=0.5em] circle (0.5em)++(0.5em,0)coordinate(t)++(-0.5em,-0.5em)coordinate(c);
\draw[-latex,thick](t)--++(1em,0)coordinate(t);
\node[draw,rectangle,anchor=west] (b1) at(t) {$\frac{C_m}{R(J_ms+f_m)}$};
\draw[-latex,thick](b1.east)--node[above,at end]{$\omega$}++(1em,0)coordinate(o)--++(1em,0)coordinate(t);
\node[draw,rectangle,anchor=west] (b4) at(t) {$\frac{1}{s}$};
\draw[-latex,thick](b4.east)--node[above,at end]{$\theta$}++(2em,0);
\draw[-latex,thick](o)|-++(-3em,-2em)coordinate(t);
\node[draw,rectangle,anchor=east] (b5) at(t) {$C_e$};
\draw[-latex,thick](b5.west)-|node[left,at end]{$-$}(c);
}
\vspace{1em}
\tikz{
\draw[-latex,thick](0,0)--node[above,at start]{$U_a$} ++(2em,0) coordinate(t);
\node[draw,rectangle,anchor=west] (b1) at(t) {$\frac{C_m/(Rf_m+C_mC_e)}{\frac{RJ_m}{Rf_m+C_mC_e}s+1}$};
\draw[-latex,thick](b1.east)--node[above,midway]{$\omega$}++(2em,0)coordinate(t);
\node[draw,rectangle,anchor=west] (b4) at(t) {$\frac{1}{s}$};
\draw[-latex,thick](b4.east)--node[above,at end]{$\theta$}++(2em,0);
}
\\
\tikz{
\draw[-latex,thick](0,0)--node[above,at start]{$U_a$} ++(2em,0) coordinate(t);
\node[draw,rectangle,anchor=west] (b1) at(t) {$\frac{K_m}{s(T_ms+1)}$};
\draw[-latex,thick](b1.east)--node[above,at end]{$\theta$}++(2em,0);
}
\end{center}
\end{frame}

\begin{frame}{\example{结构图简化}{}}
\begin{center}
\scalebox{0.8}{
\tikz{
\draw[-latex,thick](0,0)--node[above,at start]{$R$} ++(2em,0) coordinate(t);
\draw[cross] (t)[xshift=0.5em] circle (0.5em)++(0.5em,0)coordinate(t)++(-0.5em,-0.5em)coordinate(c1);
\draw[-latex,thick](t)--++(1em,0)coordinate(o1)--++(1em,0)coordinate(t);
\node[draw,rectangle,anchor=west] (b1) at(t) {$G_1$};
\draw[-latex,thick](b1.east)--++(1em,0)coordinate(t);
\draw[cross] (t)[xshift=0.5em] circle (0.5em)++(0.5em,0)coordinate(t)++(-0.5em,-0.5em)coordinate(c2)++(0,1em)coordinate(c3);
\draw[-latex,thick](t)--++(4em,0)coordinate(o2)--++(1em,0)coordinate(t);
\node[draw,rectangle,anchor=west] (b2) at(t) {$G_2$};
\draw[-latex,thick](b2.east)--++(1em,0)coordinate(t);
\draw[cross] (t)[xshift=0.5em] circle (0.5em)++(0.5em,0)coordinate(t)++(-0.5em,-0.5em)coordinate(c4);
\draw[-latex,thick](t)--++(1em,0)coordinate(o3)--++(3em,0)coordinate(o4)--++(1em,0)coordinate(t);
\node[draw,rectangle,anchor=west] (b3) at(t) {$G_3$};
\draw[-latex,thick](b3.east)--++(1em,0)coordinate(t);
\draw[cross] (t)[xshift=0.5em] circle (0.5em)++(0.5em,0)coordinate(t)++(-0.5em,0.5em)coordinate(c5);
\draw[-latex,thick](t)--++(1em,0)coordinate(o5)--node[above,at end]{$C$}++(1em,0);
\draw[-latex,thick](o1)|-++(1em,2.5em)coordinate(t);
\node[draw,rectangle,anchor=west] (b4) at(t) {$G_4$};
\draw[-latex,thick](b4.east)-|(c5);
\draw[-latex,thick](o2)|-++(-0.8em,-1.5em)coordinate(t);
\node[draw,rectangle,anchor=east] (b5) at(t) {$G_1H_1$};
\draw[-latex,thick](b5.west)-|node[left,at end]{$-$}(c2);
\draw[-latex,thick](o3)|-++(-6em,1.5em)coordinate(t);
\node[draw,rectangle,anchor=east] (b6) at(t) {$H_2$};
\draw[-latex,thick](b6.west)-|node[left,at end]{$-$}(c3);
\draw[-latex,thick](o4)|-++(-0.8em,-1.5em)coordinate(t);
\node[draw,rectangle,anchor=east] (b7) at(t) {$G_3H_3$};
\draw[-latex,thick](b7.west)-|node[left,at end]{$-$}(c4);
\draw[-latex,thick](o5)|-++(-12em,-2.5em)coordinate(t);
\node[draw,rectangle,anchor=east] (b8) at(t) {$H_4$};
\draw[-latex,thick](b8.west)-|node[left,at end]{$-$}(c1);
}}
\scalebox{0.8}{
\tikz{
\draw[-latex,thick](0,0)--node[above,at start]{$R$} ++(2em,0) coordinate(t);
\draw[cross] (t)[xshift=0.5em] circle (0.5em)++(0.5em,0)coordinate(t)++(-0.5em,-0.5em)coordinate(c1);
\draw[-latex,thick](t)--++(1em,0)coordinate(o1)--++(1em,0)coordinate(t);
\node[draw,rectangle,anchor=west] (b1) at(t) {$G_1$};
\draw[-latex,thick](b1.east)--++(1em,0)coordinate(t);
\draw[cross] (t)[xshift=0.5em] circle (0.5em)++(0.5em,0)coordinate(t)++(-0.5em,0.5em)coordinate(c2);
\draw[-latex,thick](t)--++(1em,0)coordinate(t);
\node[draw,rectangle,anchor=west] (b2) at(t) {$\frac{1}{1+G_1H_1}G_2\frac{1}{1+G_3H_3}$};
\draw[-latex,thick](b2.east)--++(1em,0)coordinate(o2)--++(1em,0)coordinate(t);
\node[draw,rectangle,anchor=west] (b3) at(t) {$G_3$};
\draw[-latex,thick](b3.east)--++(1em,0)coordinate(t);
\draw[cross] (t)[xshift=0.5em] circle (0.5em)++(0.5em,0)coordinate(t)++(-0.5em,0.5em)coordinate(c3);
\draw[-latex,thick](t)--++(1em,0)coordinate(o3)--node[above,at end]{$C$}++(1em,0);
\draw[-latex,thick](o1)|-++(1em,3em)coordinate(t);
\node[draw,rectangle,anchor=west] (b4) at(t) {$G_4$};
\draw[-latex,thick](b4.east)-|(c3);
\draw[-latex,thick](o2)|-++(-5em,2em)coordinate(t);
\node[draw,rectangle,anchor=east] (b5) at(t) {$H_2$};
\draw[-latex,thick](b5.west)-|node[left,at end]{$-$}(c2);
\draw[-latex,thick](o3)|-++(-12em,-2.1em)coordinate(t);
\node[draw,rectangle,anchor=east] (b6) at(t) {$H_4$};
\draw[-latex,thick](b6.west)-|node[left,at end]{$-$}(c1);
}}
\scalebox{0.8}{
\tikz{
\draw[-latex,thick](0,0)--node[above,at start]{$R$} ++(2em,0) coordinate(t);
\draw[cross] (t)[xshift=0.5em] circle (0.5em)++(0.5em,0)coordinate(t)++(-0.5em,-0.5em)coordinate(c1);
\draw[-latex,thick](t)--++(1em,0)coordinate(o1)--++(1em,0)coordinate(t);
\node[draw,rectangle,anchor=west] (b1) at(t) {$G_1$};
\draw[-latex,thick](b1.east)--++(1em,0)coordinate(t);
\node[draw,rectangle,anchor=west] (b2) at(t) {$\frac{G_2}{(1+G_1H_1)(1+G_3H_3)+G_2H_2}$};
\draw[-latex,thick](b2.east)--++(1em,0)coordinate(t);
\node[draw,rectangle,anchor=west] (b3) at(t) {$G_3$};
\draw[-latex,thick](b3.east)--++(1em,0)coordinate(t);
\draw[cross] (t)[xshift=0.5em] circle (0.5em)++(0.5em,0)coordinate(t)++(-0.5em,0.5em)coordinate(c2);
\draw[-latex,thick](t)--++(1em,0)coordinate(o2)--node[above,at end]{$C$}++(1em,0);
\draw[-latex,thick](o1)|-++(1em,2.1em)coordinate(t);
\node[draw,rectangle,anchor=west] (b4) at(t) {$G_4$};
\draw[-latex,thick](b4.east)-|(c2);
\draw[-latex,thick](o2)|-++(-12em,-2.1em)coordinate(t);
\node[draw,rectangle,anchor=east] (b6) at(t) {$H_4$};
\draw[-latex,thick](b6.west)-|node[left,at end]{$-$}(c1);
}}
\scalebox{0.8}{
\tikz{
\draw[-latex,thick](0,0)--node[above,at start]{$R$} ++(2em,0) coordinate(t);
\draw[cross] (t)[xshift=0.5em] circle (0.5em)++(0.5em,0)coordinate(t)++(-0.5em,-0.5em)coordinate(c);
\draw[-latex,thick](t)--++(1em,0)coordinate(t);
\node[draw,rectangle,anchor=west] (b1) at(t) {$G_4+\frac{G_1G_2G_3}{(1+G_1H_1)(1+G_3H_3)+G_2H_2}$};
\draw[-latex,thick](b1.east)--++(1em,0)coordinate(o)--node[above,at end]{$C$}++(1em,0);
\draw[-latex,thick](o)|-++(-8em,-2.1em)coordinate(t);
\node[draw,rectangle,anchor=east] (b2) at(t) {$H_4$};
\draw[-latex,thick](b2.west)-|node[left,at end]{$-$}(c);
}}
\end{center}
\end{frame}

\begin{frame}{\example{引出点移动}{}}
\begin{center}
\tikz{
\draw[-latex,thick](0,0)-- ++(2em,0) coordinate(t);
\draw[cross] (t)[xshift=0.5em] circle (0.5em)++(0.5em,0)coordinate(t)++(-0.5em,-0.5em)coordinate(c1);
\draw[-latex,thick](t)--++(1em,0)coordinate(t);
\node[draw,rectangle,anchor=west] (b1) at(t) {$G_1$};
\draw[-latex,thick](b1.east)--++(1em,0)coordinate(t);
\draw[cross] (t)[xshift=0.5em] circle (0.5em)++(0.5em,0)coordinate(t)++(-0.5em,0.5em)coordinate(c2);
\draw[-latex,thick](t)--++(1em,0)coordinate(t);
\node[draw,rectangle,anchor=west] (b2) at(t) {$G_2$};
\draw[-latex,thick](b2.east)--++(1em,0)coordinate(t);
\draw[cross] (t)[xshift=0.5em] circle (0.5em)++(0.5em,0)coordinate(t)++(-0.5em,-0.5em)coordinate(c3);
\draw[-latex,thick](t)--++(1em,0)coordinate(t);
\node[draw,rectangle,anchor=west] (b3) at(t) {$G_3$};
\draw[-latex,thick](b3.east)--++(1em,0)coordinate(o1)--++(1em,0)coordinate(t);
\node[draw,rectangle,anchor=west] (b4) at(t) {$G_4$};
\draw[-latex,thick](b4.east)--++(1em,0)coordinate(o2)--++(1em,0);
\draw[-latex,thick](o2)|-++(-2.9em,-2em)coordinate(t);
\node[draw,rectangle,anchor=east] (b5) at(t) {$H_3$};
\draw[-latex,thick](b5.west)-|node[left,at end]{$-$}(c3);
\draw[-latex,thick](o2)|-++(-10em,-3em)coordinate(t);
\node[draw,rectangle,anchor=east] (b6) at(t) {$H_1$};
\draw[-latex,thick](b6.west)-|node[left,at end]{$-$}(c1);
\draw[-latex,thick](o1)|-++(-3.5em,1.5em)coordinate(t);
\node[draw,rectangle,anchor=east] (b7) at(t) {$H_2$};
\draw[-latex,thick](b7.west)-|node[left,at end]{$-$}(c2);
}
\tikz{
\draw[-latex,thick](0,0)-- ++(2em,0) coordinate(t);
\draw[cross] (t)[xshift=0.5em] circle (0.5em)++(0.5em,0)coordinate(t)++(-0.5em,-0.5em)coordinate(c1);
\draw[-latex,thick](t)--++(1em,0)coordinate(t);
\node[draw,rectangle,anchor=west] (b1) at(t) {$G_1$};
\draw[-latex,thick](b1.east)--++(1em,0)coordinate(t);
\draw[cross] (t)[xshift=0.5em] circle (0.5em)++(0.5em,0)coordinate(t)++(-0.5em,0.5em)coordinate(c2);
\draw[-latex,thick](t)--++(1em,0)coordinate(t);
\node[draw,rectangle,anchor=west] (b2) at(t) {$G_2$};
\draw[-latex,thick](b2.east)--++(1em,0)coordinate(t);
\draw[cross] (t)[xshift=0.5em] circle (0.5em)++(0.5em,0)coordinate(t)++(-0.5em,-0.5em)coordinate(c3);
\draw[-latex,thick](t)--++(1em,0)coordinate(t);
\node[draw,rectangle,anchor=west] (b3) at(t) {$G_3$};
\draw[-latex,thick](b3.east)--++(1em,0)coordinate(o1)--++(1em,0)coordinate(t);
\node[draw,rectangle,anchor=west] (b4) at(t) {$G_4$};
\draw[-latex,thick](b4.east)--++(1em,0)coordinate(o2)--++(1em,0);
\draw[-latex,thick](o2)|-++(-2.9em,-2em)coordinate(t);
\node[draw,rectangle,anchor=east] (b5) at(t) {$H_3$};
\draw[-latex,thick](b5.west)-|node[left,at end]{$-$}(c3);
\draw[-latex,thick](o2)|-++(-10em,-3em)coordinate(t);
\node[draw,rectangle,anchor=east] (b6) at(t) {$H_1$};
\draw[-latex,thick](b6.west)-|node[left,at end]{$-$}(c1);
\draw[-latex,line width=1ex,color=cyan](o1)|-++(-3.5em,1.5em)coordinate(t);
\node[draw,rectangle,anchor=east] (b7) at(t) {$H_2$};
\draw[-latex,thick](b7.west)-|node[left,at end]{$-$}(c2);
\draw[-latex,thick](o2)|-++(-3em,1.5em)coordinate(t);
\node[draw,rectangle,anchor=east] (b8) at(t) {$\frac{1}{G_4}$};
\draw[-latex,thick](b8.west)--(b7.east);
}
\end{center}
\end{frame}


\begin{frame}{\example{综合点移动}{}}
\begin{center}
\tikz{
\draw[-latex,thick](0,0)-- ++(2em,0) coordinate(o1)--++(1em,0)coordinate(t);
\draw[cross] (t)[xshift=0.5em] circle (0.5em)++(0.5em,0)coordinate(t)++(-0.5em,-0.5em)coordinate(c1);
\draw[-latex,thick](t)--++(2em,0)coordinate(t);
\node[draw,rectangle,anchor=west] (b1) at(t) {$G_1$};
\draw[-latex,thick](b1.east)--++(2em,0)coordinate(t);
\draw[cross] (t)[xshift=0.5em] circle (0.5em)++(0.5em,0)coordinate(t)++(-0.5em,0.5em)coordinate(c2);
\draw[-latex,thick](t)--++(2em,0)coordinate(t);
\node[draw,rectangle,anchor=west] (b2) at(t) {$G_2$};
\draw[-latex,thick](b2.east)--++(2em,0)coordinate(o2)--++(2em,0)coordinate(t);
\draw[-latex,thick](o2)|-++(-1em,-3em)coordinate(t);
\node[draw,rectangle,anchor=east] (b3) at(t) {$H_1$};
\draw[-latex,thick](b3.west)-|node[left,at end]{$-$}(c1);
\draw[-latex,thick](o1)|-++(5em,3em)coordinate(t);
\node[draw,rectangle,anchor=west] (b4) at(t) {$G_3$};
\draw[-latex,thick](b4.east)-|(c2);
}
\tikz{
\draw[-latex,thick](0,0)-- ++(2em,0) coordinate(o1)-++(1em,0)coordinate(t);
\draw[cross,color=cyan] (t)[xshift=0.5em] circle (0.5em)++(0.5em,0)coordinate(t)++(-0.5em,-0.5em)coordinate(c1);
\draw[-latex,thick](t)--++(2em,0)coordinate(t);
\node[draw,rectangle,anchor=west] (b1) at(t) {$G_1$};
\draw[-latex,thick](b1.east)--++(2em,0)coordinate(t);
\draw[cross] (t)[xshift=0.5em] circle (0.5em)++(0.5em,0)coordinate(t)++(-0.5em,0.5em)coordinate(c2);
\draw[-latex,thick](t)--++(2em,0)coordinate(t);
\node[draw,rectangle,anchor=west] (b2) at(t) {$G_2$};
\draw[-latex,thick](b2.east)--++(2em,0)coordinate(o2)--++(2em,0)coordinate(t);
\draw[-latex,thick](o2)|-++(-1em,-3em)coordinate(t);
\node[draw,rectangle,anchor=east] (b3) at(t) {$H_1$};
\draw[-latex,line width=1ex,color=cyan](b3.west)-|node[left,at end]{$-$}(c1);
\draw[-latex,thick](o1)|-++(5em,3em)coordinate(t);
\node[draw,rectangle,anchor=west] (b4) at(t) {$G_3$};
\draw[-latex,thick](b4.east)-|(c2);
\draw[-latex,thick](b3.west)--++(-1em,0)coordinate(t);
\node[draw,rectangle,anchor=east] (b5) at(t) {$G_1$};
\draw[-latex,thick](b5.west)-|node[left,at end]{$-$}($ (c2)+(0,-1em)$);
}
\end{center}
\end{frame}

\begin{frame}{\example{试简化系统结构图，并求系统传递函数}{ex:simplifyTF1}}
\begin{center}
\tikz{
\draw[-latex,thick](0,0)-- node[above,at start]{$R(s)$}++(3em,0) coordinate(t);
\draw[cross] (t)[xshift=0.5em] circle (0.5em)++(0.5em,0)coordinate(t)++(-0.5em,-0.5em)coordinate(c1);
\draw[-latex,thick](t)--++(2em,0)coordinate(o1)--++(2em,0)coordinate(t);
\node[draw,rectangle,anchor=west] (b1) at(t) {$G$};
\draw[-latex,thick](b1.east)--++(2em,0)coordinate(o2)--++(2em,0)coordinate(t);
\draw[cross] (t)[xshift=0.5em] circle (0.5em)++(0.5em,0)coordinate(t)++(-0.5em,0.5em)coordinate(c2);
\draw[-latex,thick](t)--node[above,at end]{$Y(s)$}++(3em,0);
\draw[-latex,thick](o1)|-++(3em,2em)coordinate(t);
\node[draw,rectangle,anchor=west] (b2) at(t) {$H_1$};
\draw[-latex,thick](b2.east)-|(c2);
\draw[-latex,thick](o2)|-++(-3em,-2em)coordinate(t);
\node[draw,rectangle,anchor=east] (b3) at(t) {$H_2$};
\draw[-latex,thick](b3.west)-|node[left,at end]{$-$}(c1);
\draw[-latex,line width=0.5ex,color=cyan,dashed](o1)..controls($(o1)+(1em,-2em)$)and($(o2)+(-1em,-1em)$)..(o2);
\draw[color=cyan](o1)circle(5pt);
\draw[color=cyan](o2)circle(5pt);
\node[draw=orange,rectangle callout,line width=1pt, callout relative pointer={(-3em,1em)},minimum height=2em,text width=8em]at (17em,-2.5em){方法1:引出点后移};
}
\vspace{1em}\\
\tikz{
\draw[-latex,thick](0,0)-- node[above,at start]{$R(s)$}++(3em,0) coordinate(t);
\draw[cross] (t)[xshift=0.5em] circle (0.5em)++(0.5em,0)coordinate(t)++(-0.5em,-0.5em)coordinate(c1);
\draw[-latex,thick](t)--++(1em,0)coordinate(o1)--++(1em,0)coordinate(t);
\node[draw,rectangle,anchor=west] (b1) at(t) {$G$};
\draw[-latex,thick](b1.east)--++(1em,0)coordinate(o2)--++(7em,0)coordinate(t);
\draw[cross] (t)[xshift=0.5em] circle (0.5em)++(0.5em,0)coordinate(t)++(-0.5em,0.5em)coordinate(c2);
\draw[-latex,thick](t)--node[above,at end]{$Y(s)$}++(3em,0);
\draw[-latex,thick](o2)|-++(1em,2em)coordinate(t);
\node[draw,rectangle,anchor=west] (b2) at(t) {$1/G$};
\draw[-latex,thick](b2.east)--++(1em,0)coordinate(t);
\node[draw,rectangle,anchor=west] (b3) at(t) {$H_1$};
\draw[-latex,thick](b3.east)-|(c2);
\draw[-latex,thick](o2)|-++(-1em,-2em)coordinate(t);
\node[draw,rectangle,anchor=east] (b3) at(t) {$H_2$};
\draw[-latex,thick](b3.west)-|node[left,at end]{$-$}(c1);
}
$$\frac{Y(s)}{R(s)}=\frac{G+H_1}{1+GH_2}$$
\end{center}
\end{frame}

\begin{frame}{\example{试简化系统结构图，并求系统传递函数}{ex:simplifyTF2}}
\begin{center}
\tikz{
\draw[-latex,thick](0,0)-- node[above,at start]{$R(s)$}++(3em,0) coordinate(t);
\draw[cross] (t)[xshift=0.5em] circle (0.5em)++(0.5em,0)coordinate(t)++(-0.5em,-0.5em)coordinate(c1);
\draw[-latex,thick](t)--++(2em,0)coordinate(o1)--++(2em,0)coordinate(t);
\node[draw,rectangle,anchor=west] (b1) at(t) {$G$};
\draw[-latex,thick](b1.east)--++(2em,0)coordinate(o2)--++(2em,0)coordinate(t);
\draw[cross] (t)[xshift=0.5em] circle (0.5em)++(0.5em,0)coordinate(t)++(-0.5em,0.5em)coordinate(c2);
\draw[-latex,thick](t)--node[above,at end]{$Y(s)$}++(3em,0);
\draw[-latex,thick](o1)|-++(3em,2em)coordinate(t);
\node[draw,rectangle,anchor=west] (b2) at(t) {$H_1$};
\draw[-latex,thick](b2.east)-|(c2);
\draw[-latex,thick](o2)|-++(-3em,-2em)coordinate(t);
\node[draw,rectangle,anchor=east] (b3) at(t) {$H_2$};
\draw[-latex,thick](b3.west)-|node[left,at end]{$-$}(c1);
\draw[-latex,line width=0.5ex,color=cyan,dashed](o2)..controls($(o2)+(-1em,-2em)$)and($(o1)+(1em,-1em)$)..(o1);
\draw[color=cyan](o1)circle(5pt);
\draw[color=cyan](o2)circle(5pt);
\node[draw=orange,rectangle callout,line width=1pt, callout relative pointer={(3em,-1em)},minimum height=2em,text width=8em]at (0em,2.5em){方法2:引出点前移};
}
\vspace{1em}\\
\tikz{
\draw[-latex,thick](0,0)-- node[above,at start]{$R(s)$}++(3em,0) coordinate(t);
\draw[cross] (t)[xshift=0.5em] circle (0.5em)++(0.5em,0)coordinate(t)++(-0.5em,-0.5em)coordinate(c1);
\draw[-latex,thick](t)--++(6em,0)coordinate(o1)--++(2em,0)coordinate(t);
\node[draw,rectangle,anchor=west] (b1) at(t) {$G$};
\draw[-latex,thick](b1.east)--++(2em,0)coordinate(t);
\draw[cross] (t)[xshift=0.5em] circle (0.5em)++(0.5em,0)coordinate(t)++(-0.5em,0.5em)coordinate(c2);
\draw[-latex,thick](t)--node[above,at end]{$Y(s)$}++(3em,0);
\draw[-latex,thick](o1)|-++(2em,2em)coordinate(t);
\node[draw,rectangle,anchor=west] (b2) at(t) {$H_1$};
\draw[-latex,thick](b2.east)-|(c2);
\draw[-latex,thick](o1)|-++(-1em,-2em)coordinate(t);
\node[draw,rectangle,anchor=east] (b3) at(t) {$G$};
\draw[-latex,thick](b3.west)--++(-1em,0)coordinate(t);
\node[draw,rectangle,anchor=east] (b4) at(t) {$H_2$};
\draw[-latex,thick](b4.west)-|node[left,at end]{$-$}(c1);
}
$$\frac{Y(s)}{R(s)}=\frac{G+H_1}{1+GH_2}$$
\end{center}
\end{frame}


%:Lecture 4: 信号流图及梅逊公式
\lecture{信号流图及梅森公式}{Lecture 4}
\section{信号流图及梅森公式}

%\begin{frame}{主要内容}
%	\tableofcontents[currentsection]
%\end{frame}

\begin{frame}{课程回顾：控制系统的结构图及其等效变换}
\begin{block}{一、动态结构图的组成}
信号线、引出点、综合点、方框
\end{block}
\begin{block}{二、典型连接方式及等效变换}
\begin{enumerate}
\item 串联等效
\item 并联等效
\item 反馈等效
\end{enumerate}
\end{block}
\begin{block}{三、等效移动规则}
\begin{enumerate}
\item 引出点的移动
\item 综合点的移动
\end{enumerate}
\end{block}
\begin{block}{四、结构图等效变换方法}
\end{block}
\end{frame}

\subsection{信号流图的基本概念}

\begin{frame}{信号流图-拓扑结构}
	\begin{figure}
		\includegraphics[height=5cm]{../figs/C2Mobius_strip.jpeg}
		\includegraphics[height=5cm]{../figs/C2Trefoil_knot_arb.png}
		\caption{（左）莫比乌斯环和（右）三叶结（图片引用自\href{https://zh.wikipedia.org/wiki/拓扑学}{\underline{Wikipedia}} ）}
	\end{figure}
\end{frame}

%\begin{frame}{信号流图-拓扑结构}
%	\begin{columns}[T]
%		\column{.5\textwidth}
%			\animategraphics[width=.9\linewidth, autoplay=True, loop=True, palindrome=True]{12}{../figs/C2Mug_and_Torus_morph/foo-}{0}{57}\\
%			图片引用自\href{https://zh.wikipedia.org/wiki/拓扑学}{\underline{Wikipedia}} 
%		\column{.475\textwidth}
%			\animategraphics[width=.9\linewidth, autoplay=True, loop=True, palindrome=True]{12}{../figs/C2Spot_the_cow/foo-}{0}{23}
%	\end{columns}	
%\end{frame}

\begin{frame}{信号流图的基本概念}
	信号流图是一种描述系统中各信号传递关系的数学图形。只适用于线性系统。其优点在于流图增益公式实用性。信号流图由节点和支路组成。
	\begin{center}
		\tikz[node distance=4em, bend angle=60]{
		\node[var] (x1) [label=left:$x_1$]{};
		\node[var] (x2) [right of=x1, label=above left:$x_2$]{};
		\node[var] (x3) [right of=x2, label=above right:$x_3$]{};
		\node[var] (x4) [right of=x3, label=right:$x_4$]{};
		\path[con] (x1) -- node[above]{a}(x2);
		\path[con] (x2) -- node[above]{b}(x3);
		\path[con] (x3) -- node[above]{1}(x4);
		\path[con] (x3) edge [bend right]node[above]{c}(x2);
		}
	\end{center}
	\begin{description}
		\item[节点\tikz{\node[var]{};}：]表示系统中的变量。
		\item[支路\tikz{\path[con](0,0)--(2em,0);}：] 表示变量之间的传输关系。
	\end{description}
\end{frame}

\begin{frame}{信流图的基本术语}
	\begin{description}
		\onslide<1->{\item[源点：]\ }\onslide<2->{只有流出支路的节点。对应于系统的输入信号，或称为输入节点。}
		\onslide<1->{\item[陷点：]\ }\onslide<3->{只有输入支路的节点。对应于系统的输出信号，或称为输出节点。}
		\onslide<1->{\item[混合节点：]\ }\onslide<4->{既有输入也有输出的节点。}
	\end{description}
	\tikzstyle{call} = [draw=orange, rectangle callout, line width=1pt, minimum height=2em, text width=4em]
\begin{center}
	\tikz[node distance=4em, bend angle=60]{
	\node[var](x1)[label=left:$x_1$]{};
	\node[var](x2)[right of=x1, label=above left:$x_2$]{};
	\node[var](x3)[right of=x2, label=above right:$x_3$]{};
	\node[var](x4)[above of=x3, label=left:$x_4$]{};
	\node[var](x5)[right of=x3, label=right:$x_5$]{};
	\path[con](x1) -- node[above]{a}(x2);
	\path[con](x2) -- node[above]{b}(x3);
	\path[con](x3) -- node[above]{e}(x5);
	\path[con](x4) -- node[right]{d}(x3);
	\path[con](x3)edge[bend left]node[above]{c}(x2);
	\onslide<2->{\node[call, callout relative pointer={(0em,1em)}] at (0em,-3em){输入节点（源点）};
	\node[call, callout relative pointer={(-1em,0em)}] at (12em,4em){输入节点（源点）};}
	\onslide<3->{\node[call, callout relative pointer={(0em,1em)}] at (12em,-3em){输出节点（陷点）};}
	\onslide<4->{\node[call, callout relative pointer={(0em,-1em)}] at (4em,2.5em){混合节点};}
	}
\end{center}
\end{frame}

\begin{frame}{信流图的基本术语}
	\begin{description}
		\onslide<1->{\item[通路：]\ }\onslide<2->{从某一节点开始沿支路箭头方向到另一节点（或同一节点）构成的路径。}
		\onslide<1->{\item[回路：]\ }\onslide<3->{如果通路的终点就是起点，并且与任何其他节点相交不多于一次。}
		\onslide<1->{\item[回环增益：]\ }\onslide<4->{回环中各支路增益的乘积。}
		\onslide<1->{\item[前向通路：]\ }\onslide<5->{从源点开始并终止于陷点，且与其他节点相交不多于一次的通道。该通路各支路增益乘积称为前向通路增益。}
		\onslide<1->{\item[不接触回路：]\ }\onslide<6->{各回路之间没有任何公共节点。反之称为接触回路 。}
	\end{description}
	\begin{center}
		\tikz[node distance=3em,bend angle=60]{
		\node[var](R)[label=left:$R$]{};
		\foreach \x/\y/\z/\n in {e/R/above/e, a1/e/above/a_1, a2/a1/below/a_2, a3/a2/above/a_3, a4/a3/above/a_4, a5/a4/below/a_5, a6/a5/above/a_6, a7/a6/above/a_7, C/a7/right/C}{
			\node[var](\x)[right of=\y, label=\z:$\n$]{};
		}	
		\foreach \x/\y/\z in {R/1/e, e/G_1/a1, a1/G_2/a2, a2/G_3/a3, a3/1/a4, a4/G_4/a5, a5/G_5/a6, a6/G_6/a7, a7/1/C}{
			\path[con](\x)--node[above]{$\y$}(\z);
		}
		\path[con](a5)edge[bend right]node[below]{$-H_4$}node[above, near start]{\rm II}(a2);
		\path[con](a7)to[bend left]node[above]{$-H_1$}node[above, near end]{\rm IV}(e);
		\path[con](a3)edge[bend left]node[below]{$-H_2$}node[above, near start]{\rm I}(a1);
		\path[con](a6)edge[bend left]node[below]{$-H_3$}node[above, near start]{\rm III}(a4);
		}
	\end{center}
\end{frame}

\subsection{信号流图的绘制}
\begin{frame}{信号流图的绘制}
	\begin{block}{由结构图绘制信流图}
		\begin{align*}
			\text{\alert{结构图}} &	&\text{\alert{信号流图}}\\
			\text{输入信号} &\hspace{4em}\longrightarrow&\text{源节点}\\
			\text{输出信号} &\hspace{4em}\longrightarrow&\text{陷节点}\\
			\text{比较点，引出点} &\hspace{4em}\longrightarrow&\text{混合节点 }\\
			\text{环节} &\hspace{4em}\longrightarrow&\text{支路}\\
			\text{环节传递函数} &\hspace{4em}\longrightarrow&\text{支路增益 }
		\end{align*}
	\end{block}
\end{frame}

\begin{frame}{信号流图的绘制}
	\begin{block}{基本步骤：}
		\begin{enumerate}
			\item 在结构图的信号线上，标出各变量对应节点名称。若比较点之后又有引出点，只需设一个节点。
			\item 所有输入为源点，输出为陷点，可以通过引入单位增益支路实现。
			\item 将各节点按原来顺序排列；
			\item 连接各支路，注意支路增益符号。 
		\end{enumerate}
	\end{block}
\end{frame}

\begin{frame}{由结构图绘制信号流图}
\begin{center}
\tikz{
        \draw[-latex,thick](0,0)-- node[above,at start]{$r$}++(1em,0) coordinate(t);
        \draw[cross] (t)[xshift=0.5em] circle (0.5em)++(0.5em,0)coordinate(t)++(-0.5em,-0.5em)coordinate(c1);
        \draw[-latex,thick](t)--node[above,midway]{$e$}++(1em,0)coordinate(t);
        \node[draw,rectangle,anchor=west] (b1) at(t) {$G_1$};
        \draw[-latex,thick](b1.east)--++(1em,0)coordinate(t);
        \draw[cross] (t)[xshift=0.5em] circle (0.5em)++(0.5em,0)coordinate(t)++(-0.5em,-0.5em)coordinate(c2);
        \draw[-latex,thick](t)--node[above,midway]{$a_1$}++(1em,0)coordinate(t);
        \node[draw,rectangle,anchor=west] (b2) at(t) {$G_2$};
        \draw[-latex,thick](b2.east)--++(1em,0)coordinate(t);
        \draw[cross] (t)[xshift=0.5em] circle (0.5em)++(0.5em,0)coordinate(t)++(-0.5em,0.5em)coordinate(c3);
        \draw[-latex,thick](t)--node[below,midway]{$a_2$}++(1em,0)coordinate(t);
        \node[draw,rectangle,anchor=west] (b3) at(t) {$G_3$};
        \draw[-latex,thick](b3.east)--node[above,at end]{$a_3$}++(1em,0)coordinate(o1)--++(1em,0)coordinate(t);
        \draw[cross] (t)[xshift=0.5em] circle (0.5em)++(0.5em,0)coordinate(t)++(-0.5em,-0.5em)coordinate(c4);
        \draw[-latex,thick](t)--node[above,midway]{$a_4$}++(1em,0)coordinate(t);
        \node[draw,rectangle,anchor=west] (b4) at(t) {$G_4$};
        \draw[-latex,thick](b4.east)--++(1em,0)coordinate(o2)--node[above,midway]{$a_5$}++(1em,0)coordinate(t);
        \node[draw,rectangle,anchor=west] (b5) at(t) {$G_5$};
        \draw[-latex,thick](b5.east)--++(1em,0)coordinate(o3)--node[above,at start]{$a_6$}++(1em,0)coordinate(t);
        \node[draw,rectangle,anchor=west] (b6) at(t) {$G_6$};
        \draw[-latex,thick](b6.east)--++(1em,0)coordinate(o4)--node[above,near start]{$a_7=c$}++(1em,0);
        \draw[-latex,thick](o1)|-++(-3em,-2em)coordinate(t);
        \node[draw,rectangle,anchor=east] (h2) at(t) {$H_2$};
        \draw[-latex,thick](h2.west)-|node[right,at end]{$-$}(c2);
        \draw[-latex,thick](o3)|-++(-3em,-2em)coordinate(t);
        \node[draw,rectangle,anchor=east] (h3) at(t) {$H_3$};
        \draw[-latex,thick](h3.west)-|node[right,at end]{$-$}(c4);
        \draw[-latex,thick](o2)|-++(-4em,2em)coordinate(t);
        \node[draw,rectangle,anchor=east] (h4) at(t) {$H_4$};
        \draw[-latex,thick](h4.west)-|node[right,at end]{$-$}(c3);
        \draw[-latex,thick](o4)|-++(-11em,-3em)coordinate(t);
        \node[draw,rectangle,anchor=east] (h1) at(t) {$H_1$};
        \draw[-latex,thick](h1.west)-|node[right,at end]{$-$}(c1);
} 
\vspace{1em}
\tikz[node distance=3em,bend angle=60]{
\onslide<2->{
\node[var](R)[label=left:$R$]{};
\node[var](e)[right of=R,label=above:$e$]{};
\node[var](a1)[right of=e,label=above:$a_1$]{};
\node[var](a2)[right of=a1,label=below:$a_2$]{};
\node[var](a3)[right of=a2,label=above:$a_3$]{};
\node[var](a4)[right of=a3,label=above:$a_4$]{};
\node[var](a5)[right of=a4,label=below:$a_5$]{};
\node[var](a6)[right of=a5,label=above:$a_6$]{};
\node[var](a7)[right of=a6,label=above:$a_7$]{};
\node[var](C)[right of=a7,label=right:$C$]{};}
\onslide<3->{
\path[con](R)--node[above]{$1$}(e);
\path[con](e)--node[above]{$G_1$}(a1);
\path[con](a1)--node[above]{$G_2$}(a2);
\path[con](a2)--node[above]{$G_3$}(a3);
\path[con](a3)--node[above]{$1$}(a4);
\path[con](a4)--node[above]{$G_4$}(a5);
\path[con](a5)--node[above]{$G_5$}(a6);
\path[con](a6)--node[above]{$G_6$}(a7);
\path[con](a7)--node[above]{$1$}(C);
\path[con](a5)edge[bend right]node[below]{$-H_4$}(a2);
\path[con](a7)to[bend left]node[above]{$-H_1$}(e);
\path[con](a3)edge[bend left]node[below]{$-H_2$}(a1);
\path[con](a6)edge[bend left]node[below]{$-H_3$}(a4);}
}
\end{center}
\end{frame}

\begin{frame}{由结构图绘制信号流图}
\begin{center}
\tikzstyle{var}=[draw,color=orange,thick,circle,minimum size=0.2em,inner sep=2pt]
\tikzstyle{con}=[draw,color=cyan,thick,-latex]
\tikz{
\draw[-latex,thick](0,0)-- node[above,at start]{$r$}++(1em,0) coordinate(t);
\draw[cross] (t)[xshift=0.5em] circle (0.5em)++(0.5em,0)coordinate(t)++(-0.5em,-0.5em)coordinate(c1);
\draw[-latex,thick](t)--node[above,midway]{$e$}++(1em,0)coordinate(t);
\node[draw,rectangle,anchor=west] (b1) at(t) {$K$};
\draw[-latex,thick](b1.east)--++(1em,0)coordinate(t);
\draw[cross] (t)[xshift=0.5em] circle (0.5em)++(0.5em,0)coordinate(t)++(-0.5em,0.5em)coordinate(c2);
\draw[-latex,thick](t)--node[below,midway]{$a_1$}++(1em,0)coordinate(t);
\node[draw,rectangle,anchor=west] (b2) at(t) {$0.5$};
\draw[-latex,thick](b2.east)--++(1em,0)coordinate(t);
\draw[cross] (t)[xshift=0.5em] circle (0.5em)++(0.5em,0)coordinate(t)++(-0.5em,-0.5em)coordinate(c3);
\draw[-latex,thick](t)--node[above,midway]{$a_2$}++(1em,0)coordinate(t);
\node[draw,rectangle,anchor=west] (b3) at(t) {$\frac{1}{s+1}$};
\draw[-latex,thick](b3.east)--node[below,at end]{$a_3$}++(1em,0)coordinate(o1)--++(1em,0)coordinate(t);
\node[draw,rectangle,anchor=west] (b4) at(t) {$\frac{1}{s^2}$};
\draw[-latex,thick](b4.east)--++(1em,0)coordinate(o2)--node[above,midway]{$a_4=c$}++(1em,0);
\draw[-latex,thick](o1)|-++(-3em,2em)coordinate(t);
\node[draw,rectangle,anchor=east] (h1) at(t) {$5$};
\draw[-latex,thick](h1.west)-|node[right,at end]{$-$}(c2);
\draw[-latex,thick](o2)|-++(-3em,-2em)coordinate(t);
\node[draw,rectangle,anchor=east] (h2) at(t) {$s$};
\draw[-latex,thick](h2.west)-|node[right,at end]{$-$}(c3);
\draw[-latex,thick](o2)--++(0,-3em)-|node[right,at end]{$-$}(c1);
} 
\vspace{1em}
\tikz[node distance=4em,bend angle=60]{
\onslide<2->{
\node[var](r)[label=left:$r$]{};
\node[var](e)[right of=r,label=above:$e$]{};
\node[var](a1)[right of=e,label=above:$a_1$]{};
\node[var](a2)[right of=a1,label=above:$a_2$]{};
\node[var](a3)[right of=a2,label=below:$a_3$]{};
\node[var](a4)[right of=a3,label=above:$a_4$]{};
\node[var](c)[right of=a4,label=right:$c$]{};}
\onslide<3->{
\path[con](r)--node[above]{$1$}(e);
\path[con](e)--node[above]{$K$}(a1);
\path[con](a1)--node[above]{$0.5$}(a2);
\path[con](a2)--node[above]{$\frac{1}{s+1}$}(a3);
\path[con](a3)--node[above]{$\frac{1}{s^2}$}(a4);
\path[con](a4)--node[above]{$1$}(c);
\path[con](a3)edge[bend right]node[above]{$-5$}(a1);
\path[con](a4)edge[bend left]node[below]{$-s$}(a2);
\path[con](a4)edge[bend left]node[above]{$-1$}(e);}
}
\end{center}
\end{frame}

\begin{frame}{由方程组绘制信号流图}
首先按照节点的次序绘出各节点，然后根据各方程式绘制各支路。当所有方程式的信号流图绘制完毕后，即得系统的信号流图。
\begin{eqnarray*}
&x_1=x_r-gx_c\quad x_2=ax_1-fx_4\quad x_3=bx_2-ex_c\\
&x_4=cx_3\quad x_c=dx_4
\end{eqnarray*}
\begin{center}
\tikz[node distance=4em,bend angle=60]{
\onslide<2->{
\node[var](xr)[label=left:$x_r$]{};
\node[var](x1)[right of=xr,label=above:$x_1$]{};
\node[var](x2)[right of=x1,label=below:$x_2$]{};
\node[var](x3)[right of=x2,label=above:$x_3$]{};
\node[var](x4)[right of=x3,label=above:$x_4$]{};
\node[var](xc)[right of=x4,label=right:$x_c$]{};}
\onslide<3->{
\path[con](xr)--node[above]{$1$}(x1);
\path[con](x1)--node[above]{$a$}(x2);
\path[con](x2)--node[above]{$b$}(x3);
\path[con](x3)--node[above]{$c$}(x4);
\path[con](x4)--node[above]{$d$}(xc);
\path[con](xc)edge[bend left]node[above]{$-e$}(x3);
\path[con](x4)edge[bend right]node[above]{$-f$}(x2);
\path[con](xc)edge[bend left]node[above]{$-g$}(x1);}
}
\end{center} 
\end{frame}

\subsection{梅森（Mason）增益公式}

\begin{frame}{如何求解信号的增益？}
%	\begin{figure}
	\begin{tikzpicture}[x=1em, y=1em, node distance=4em, bend angle=60]
		\node[var] (xr) [label=left:$x_r$] {};
		\node[var](x1)[right of=xr, label=above:$x_1$]{};
		\node[var](xc)[right of=x1, label=right:$x_c$]{};
		\path[con](xr)--node[above] {$a$} (x1);
		\path[con](x1)--node[above] {$b$} (xc);
		\begin{scope}[yshift=-7em]
			\node[var] (xr) [label=left:$x_r$] {};
			\node[var] (xc)[right of=xr, label=right:$x_c$] {};
			\path[con] (xr) edge [bend left] node [above] {$a$} (xc);
			\path[con] (xr) edge [bend right] node [below] {$b$} (xc);
		\end{scope}
		\begin{scope}[yshift=-14em]
			\node[var] (xr) [label=left:$x_r$] {};
			\node[var](x1)[right of=xr, label=above:$x_1$]{};
			\node[var](xc)[right of=x1, label=right:$x_c$]{};
			\path[con] (xr) -- node [above] {$1$} (x1);
			\path[con] (x1) -- node [above] {$a$} (xc);
			\path[con] (xc) edge [bend left] node [below] {$b$} (x1);
		\end{scope}
	\end{tikzpicture}
%	\end{figure}
\end{frame}

\begin{frame}{梅森（Mason）生平}
\begin{columns}[T]
	\column{.5\textwidth}
	\begin{block}{Samuel Jefferson Mason (1921–1974)}
	\begin{itemize}
	\item Professor and Associate Director of the Research Laboratory of Electronics, MIT.
	\item Leader of the Cognitive Information Processing Unit at Research Laboratory of Electronics, MIT.
	\item 博士论文：On the logic of feedback，信号流图发明者（一说为 Claude Shannon）
	\item 为盲人设计了OCR及语音合成系统
	\end{itemize}
	\end{block}
	\column{.5\textwidth}
	\includegraphics[height=7cm]{../figs/C2Mason.jpg}
\end{columns}
\end{frame}

\begin{frame}{梅森（Mason）增益公式}
$$p=\frac{1}{\Delta}\sum_{k=1}^np_k\Delta_k$$
\begin{description}
\item[$p$]总增益
\item[$p_k$]第$k$条前向通路的通路增益
\item[$\Delta$]信号流图的特征式，即
$$\Delta=1-\sum_aL_a+\sum_{bc}L_bL_c-\sum_{def}L_dL_eL_f+\cdots$$
\item[$\sum_aL_a$]所有回路增益之和
\item[$\sum_{bc}L_bL_c$]每两互不接触回路增益乘积之和
\item[$\sum_{def}L_dL_eL_f$]每三个互不接触回路增益乘积之和
\item[$\Delta_k$]在$\Delta$中除去与第$k$条前向通道相接触的回路后的特征式，称为前向通道特征式的余子式。
\end{description}
\end{frame}

\begin{frame}{由信号流图计算传递函数}
\begin{center}
\tikzstyle{var}=[draw,color=orange,thick,circle,minimum size=0.2em,inner sep=2pt]
\tikzstyle{con}=[draw,color=cyan,thick,-latex]
\tikz[node distance=4em,bend angle=60]{
\node[var](r)[label=left:$R(s)$]{};
\node[var](x1)[right of=r]{};
\node[var](x2)[right of=x1]{};
\node[var](x3)[right of=x2]{};
\node[var](x4)[right of=x3]{};
\node[var](c)[right of=x4,label=right:$C(s)$]{};
\path[con](r)--node[above]{$1$}(x1);
\path[con](x1)--node[above]{$a$}(x2);
\path[con](x2)--node[above]{$b$}(x3);
\path[con](x3)--node[above]{$c$}(x4);
\path[con](x4)--node[above]{$d$}(c);
\path[con](x1)edge[bend left]node[above]{$e$}(x4);
\path[con](x2)edge[bend left]node[below]{$f$}(x1);
\path[con](x3)edge[bend right]node[above]{$g$}(x2);
\path[con](x4)edge[bend left]node[below]{$h$}(x3);
}\\
\begin{columns}[T]
	\column{.5\textwidth}
		\begin{description}			
			\item<2->[回路增益：]$af$，$bg$，$ch$，$ehgf$
			\item<3->[每两互不接触回路：]$afch$
			\item<5->[前向通路增益：]$p_1=abcd$，$p_2=ed$
		\end{description}
	\column{.5\textwidth}
		\onslide<4->{
		$$\Delta=1-af-bg-ch-ehgf+afch$$
		}
		\onslide<6->{
		$$\Delta_1=1\quad\Delta_2=1-bg$$
		}
\end{columns}
\onslide<7->{
\alert{四个单独回路，两个回路互不接触}\\
\alert{前向通路两条}}
\end{center} 
\onslide<8->{
$$\frac{C(s)}{R(s)}=\frac{abcd+ed(1-bg)}{1-af-bg-ch-ehgf+afch}$$}
\end{frame}

\addtocounter{example}{1}
\begin{frame}{\examplec{求$C(s)/R(s)$：A}{ex:CbyR1}}
\begin{figure}
	\begin{tikzpicture}[x=1em, y=1em]
		\draw[-latex,thick](0,0)-- node[above,at start]{$R$}++(1,0) coordinate(t);
		\node[cross] (s1) at (t)[xshift=0.5em] {};
		\draw[-latex,thick](s1)--node[above,midway]{$e$}++(2,0)coordinate(t);
		\node[draw,rectangle,anchor=west] (b1) at(t) {$G_1$};
		\draw[-latex,thick](b1.east)--++(1,0)coordinate(t);
		\node[cross] (s2) at (t)[xshift=0.5em] {};
		\draw[-latex,thick](s2)--node[above,midway]{$a_1$}++(2,0) coordinate(t);
		\node[draw,rectangle,anchor=west] (b2) at(t) {$G_2$};
		\draw[-latex,thick](b2.east)--++(1em,0)coordinate(t);
		\node[cross] (s3) at (t) [xshift=0.5em] {};
		\draw[-latex,thick](s3)--node[below,midway]{$a_2$}++(2,0)coordinate(t);
		\node[draw,rectangle,anchor=west] (b3) at(t) {$G_3$};
		\draw[-latex,thick](b3.east)--++(.5em,0)coordinate(o1)--node[above,at start]{$a_3$}++(1,0)coordinate(t);
		\node[cross] (s4) at (t) [xshift=0.5em] {};
		\draw[-latex,thick](s4)--node[above,midway]{$a_4$}++(2,0)coordinate(t);
		\node[draw,rectangle,anchor=west] (b4) at(t) {$G_4$};
		\draw[-latex,thick](b4.east)--++(.5em,0)coordinate(o2)--node[above,midway]{$a_5$}++(1em,0)coordinate(t);
		\node[draw,rectangle,anchor=west] (b5) at(t) {$G_5$};
		\draw[-latex,thick](b5.east)--++(.5em,0)coordinate(o3)--node[above,at start]{$a_6$}++(1em,0)coordinate(t);
		\node[draw,rectangle,anchor=west] (b6) at(t) {$G_6$};
		\draw[-latex,thick](b6.east)--++(.5em,0)coordinate(o4)--node[above,near start]{$C$}++(1em,0);
		\draw[-latex,thick](o1)|-++(-3em,-2em)coordinate(t);
		\node[draw,rectangle,anchor=east] (h2) at(t) {$H_2$};
		\draw[-latex,thick](h2.west)-|node[right,at end]{$-$}(s2);
		\draw[-latex,thick](o3)|-++(-3em,-2em)coordinate(t);
		\node[draw,rectangle,anchor=east] (h3) at(t) {$H_3$};
		\draw[-latex,thick](h3.west)-|node[right,at end]{$-$}(s4);
		\draw[-latex,thick](o2)|-++(-4em,2em)coordinate(t);
		\node[draw,rectangle,anchor=east] (h4) at(t) {$H_4$};
		\draw[-latex,thick](h4.west)-|node[right,at end]{$-$}(s3);
		\draw[-latex,thick](o4)|-++(-11em,-3em)coordinate(t);
		\node[draw,rectangle,anchor=east] (h1) at(t) {$H_1$};
		\draw[-latex,thick](h1.west)-|node[right,at end]{$-$}(s1);
		\onslide<2,6>{\path[flow] (b2) -- (b3) --(o1) |- (h2) -| (s2.center) -- (b2);}
		\onslide<3,6>{\path[flow] (b4) -- (b5) --(o3) |- (h3) -| (s4.center) -- (b4);}
		\onslide<4>{\path[flow] (h4) -| (s3.center) --(b3) -- (b4) -- (o2) |- (h4);}
		\onslide<5>{\path[flow] (b1) -- (b2) -- (b3) -- (b4) -- (b5) -- (b6) --(o4) |- (h1) -| (s1.center) -- (b1);}
		\onslide<8>{\path[flow] (s1) -- (b1) -- (b2) -- (b3) -- (b4) -- (b5) -- (b6) -- (o4);}
	\end{tikzpicture}
\end{figure}
\begin{align*}
	\onslide<2->{\Delta&=1-[-G_2G_3H_2}\onslide<3->{-G_4G_5H_3}\onslide<4->{-G_3G_4H_4}\onslide<5->{-G_1G_2G_3G_4G_5G_6H_1]}\\
	\onslide<6->{&+(-G_2G_3H_2)(-G_4G_5H_3)}\\
	\onslide<7->{&=1+G_2G_3H_2+G_4G_5H_3+G_3G_4H_4\\
	&+G_1G_2G_3G_4G_5G_6H_1+G_2G_3G_4G_5H_2H_3}\\
	\onslide<8->{P_1&=G_1G_2G_3G_4G_5G_6~~~~\Delta_1=1}\\
	\onslide<9->{\Phi&=\frac{G_1G_2G_3G_4G_5G_6}{1+G_2G_3H_2+G_4G_5H_3+G_3G_4H_4+G_{1-6}H_1+G_{2-5}H_2H_3}}
\end{align*}
\end{frame}

\addtocounter{example}{1}
\begin{frame}{\examplec{求$C(s)/R(s)$：B}{ex:CbyR2}}
\begin{center}
	\begin{tikzpicture}[x=1em, y=1em]
		\draw[-latex,thick](0,0)-- node[above,at start]{$R$}++(1em,0) coordinate(t);
		\node[cross] (s1) at (t)[xshift=0.5em] {};
		\draw[-latex,thick](s1)--++(2,0)coordinate(t);
		\node[cross] (s2) at (t)[xshift=0.5em] {}; 
		\draw[-latex,thick](s2)--++(2,0)coordinate(t);
		\node[draw,rectangle,anchor=west] (b1) at(t) {$G_1$};
		\draw[-latex,thick](b1.east)--++(2,0)coordinate(t);
		\node[cross] (s3) at (t)[xshift=0.5em] {};
		\draw[-latex,thick](s3)--++(1em,0)coordinate(o1)--++(1em,0)coordinate(t);
		\node[draw,rectangle,anchor=west] (b2) at (t) {$G_2$};
		\draw[-latex,thick](b2.east)--++(2,0)coordinate(o2)--++(1em,0)coordinate(t);
		\node[draw,rectangle,anchor=west] (b3) at (t) {$G_3$};
		\draw[-latex,thick](b3.east)--++(2,0)coordinate(t);
		\node[cross] (s4) at (t)[xshift=0.5em] {};
		\draw[-latex,thick](s4)--++(1em,0)coordinate(o3)--node[above,at start]{$C$}++(1em,0)coordinate(t);
		\draw[-latex,thick](o2)|-++(-3em,-3em)coordinate(t);
		\node[draw,rectangle,anchor=east] (h1) at(t) {$H_1$};
		\draw[-latex,thick](h1.west)-|node[right,at end]{$-$}(s2);
		\draw[-latex,thick](o3)|-++(-3em,-2em)coordinate(t);
		\node[draw,rectangle,anchor=east] (h2) at(t) {$H_2$};
		\draw[-latex,thick](h2.west)-|node[right,at end]{$-$}(s3);
		\draw[-latex,thick](o3)--++(0,-4em)-|node[right,at end]{$-$}(s1);
		\draw[-latex,thick](o1)|-++(3em,3em)coordinate(t);
		\node[draw,rectangle,anchor=west] (b4) at(t) {$G_4$};
		\draw[-latex,thick](b4.east)-|(s4);
		\onslide<2>{\path[flow] (b1) -- (o2) --++(0,-3) -| (s2) -- (b1);}
		\onslide<3>{\draw[flow] (b2) -- (b3) -- (o3) --++(0,-2) -| (s3) -- (b2);}
		\onslide<4>{\draw[flow] (b1) -- (b2) -- (b3) -- (o3) --++(0,-4) -| (s1) -- (b1);}
		\onslide<5>{\draw[flow] (b4) -| (s4) -- (o3) --++(0,-2) -| (s3) -- (o1) |- (b4);}
		\onslide<6>{\draw[flow] (b1) -- (o1) |- (b4) -| (s4) -- (o3) --++(0,-4) -| (s1) -- (b1);}
		\onslide<8>{\draw[flow] (s1) -- (b1) -- (b2) -- (b3) -- (o3);}
		\onslide<9>{\draw[flow] (s1) -- (b1) -- (o1) |- (b4) -| (s4.center) -- (o3);}
	\end{tikzpicture}
\end{center}
\begin{align*}
\onslide<2->{
\Delta&=1-[-G_1G_2H_1\onslide<3->{-G_2G_3H_2}\onslide<4->{-G_1G_2G_3}\onslide<5->{-G_4H_2}\onslide<6->{-G_1G_4]}\\}
\onslide<7->{
&=1+G_1G_2H_1+G_2G_3H_2+G_1G_2G_3+G_4H_2+G_1G_4\\}
\onslide<8->{
P_1&=G_1G_2G_3\quad \quad \Delta_1=1\\}
\onslide<9->{P_2&=G_1G_4\quad \quad \Delta_2=1\\}
\onslide<10->{
\Phi&=\frac{G_1G_2G_3+G_1G_4}{1+G_1G_2H_1+G_2G_3H_2+G_1G_2G_3+G_4H_2+G_1G_4}}
\end{align*}
\end{frame}

\addtocounter{example}{1}
\begin{frame}{\examplec{求$C(s)/R(s)$：C}{ex:CbyR3}}
\begin{center}
	\begin{tikzpicture}[x=1em, y=1em]
		\draw[-latex,thick](0,0)-- node[above,at start]{$R$}++(1em,0) coordinate(t);
		\node[cross] (s1) at (t)[xshift=0.5em] {};%circle (0.5em)++(0.5em,0)coordinate(t)++(-0.5em,-0.5em)coordinate(c1);
		\draw[-latex,thick](s1)--++(1.5,0)coordinate(t);
		\node[draw,rectangle,anchor=west] (b1) at(t) {$\frac{1}{R}$};
		\draw[-latex,thick](b1.east)--++(1em,0)coordinate(t);
		\node[cross] (s2) at (t) [xshift=0.5em] {};%circle (0.5em)++(0.5em,0)coordinate(t)++(-0.5em,0.5em)coordinate(c2);
		\draw[-latex,thick](s2)--++(1.5,0)coordinate(t);
		\node[draw,rectangle,anchor=west] (b2) at(t) {$\frac{1}{Cs}$};
		\draw[-latex,thick](b2.east)--++(.5em,0)coordinate(o1)--++(1,0)coordinate(t);
		\node[cross] (s3) at (t) [xshift=0.5em] {};%circle (0.5em)++(0.5em,0)coordinate(t)++(-0.5em,-0.5em)coordinate(c3);
		\draw[-latex,thick](s3)--++(1.5,0)coordinate(t);
		\node[draw,rectangle,anchor=west] (b3) at(t) {$\frac{1}{R}$};
		\draw[-latex,thick](b3.east)--++(.5em,0)coordinate(o2)--++(1,0)coordinate(t);
		\node[cross] (s4) at (t)[xshift=0.5em] {};%circle (0.5em)++(0.5em,0)coordinate(t)++(-0.5em,0.5em)coordinate(c4);
		\draw[-latex,thick](s4)--++(1.5,0)coordinate(t);
		\node[draw,rectangle,anchor=west] (b4) at(t) {$\frac{1}{Cs}$};
		\draw[-latex,thick](b4.east)--++(.5em,0)coordinate(o3)--++(1,0)coordinate(t);
		\node[cross] (s5) at (t)[xshift=0.5em] {};%circle (0.5em)++(0.5em,0)coordinate(t)++(-0.5em,-0.5em)coordinate(c5);
		\draw[-latex,thick](s5)--++(1.5,0)coordinate(t);
		\node[draw,rectangle,anchor=west] (b5) at(t) {$\frac{1}{R}$};
		\draw[-latex,thick](b5.east)--++(.5em,0)coordinate(o4)--++(1,0)coordinate(t);
		\node[draw,rectangle,anchor=west] (b6) at(t) {$\frac{1}{Cs}$};
		\draw[-latex,thick](b6.east)--++(.5em,0)coordinate(o5)--node[above,at start]{$C$}++(1em,0);
		\draw[-latex,thick](o1)--++(0,-2em)-|node[above,near start]{\rm I}node[right,at end]{$-$}(s1);
		\draw[-latex,thick](o2)--++(0,2em)-|node[above,near start]{\rm II}node[right,at end]{$-$}(s2);
		\draw[-latex,thick](o3)--++(0,-2em)-|node[above,near start]{\rm III}node[right,at end]{$-$}(s3);
		\draw[-latex,thick](o4)--++(0,2em)-|node[above,near start]{\rm IV}node[right,at end]{$-$}(s4);
		\draw[-latex,thick](o5)--++(0,-2em)-|node[above,near start]{\rm V}node[right,at end]{$-$}(s5);
	\end{tikzpicture}\\
\onslide<2->{回路：{\rm I，II，III，IV，V}\\}
\onslide<3->{互不接触回路：}\onslide<3->{{\rm I-III，I-V，III-V，I-IV，II-V，II-IV}}
\end{center}
\begin{align*}
\onslide<4->{
\Delta&=1-5\cdot\frac{-1}{RCs}+6\cdot\frac{1}{(RCs)^2}-\frac{-1}{(RCs)^3}\\}
\onslide<5->{
p_1&=\frac{1}{(RCs)^3}\quad \quad \Delta_1=1\\}
\onslide<6->{
\Phi&=\frac{1}{1+(RCs)^3+5(RCs)^2+6(RCs)}}
\end{align*}
\end{frame}

\addtocounter{example}{1}
\begin{frame}{\examplec{求$C(s)/R(s)$：D}{CbyR4}}
\begin{center}
	\begin{tikzpicture}[x=1em, y=1em]
		\draw[-latex,thick](0,0)--node[above,at start]{$R$}++(1em,0) coordinate(t);
		\node[cross] (s1) at (t)[xshift=0.5em] {};
		\draw[-latex,thick](s1)--++(1,0)coordinate(t);
		\node[cross] (s2) at (t)[xshift=0.5em] {};
		\draw[-latex,thick](s2)--++(.5,0)coordinate(o1)--++(1em,0)coordinate(t);
		\node[draw,rectangle,anchor=west] (b1) at(t) {$G_1$};
		\draw[-latex,thick](b1.east)--++(1em,0)coordinate(t);
		\node[cross] (s3) at (t)[xshift=0.5em] {};
		\draw[-latex,thick](s3)--++(4em,0)coordinate(o2)--++(1em,0)coordinate(t);
		\node[draw,rectangle,anchor=west] (b2) at(t) {$G_2$};
		\draw[-latex,thick](b2.east)--++(1em,0)coordinate(t);
		\node[cross] (s4) at (t)[xshift=0.5em] {};
		\draw[-latex,thick](s4)--++(1em,0)coordinate(o3)--++(1em,0)coordinate(o4)--node[above,at start]{$C$}++(1em,0)coordinate(t);
		\draw[-latex,thick](o4)--++(0,-3)-|node[right,at end]{$-$}(s1);
		\draw[-latex,thick](o2)--++(0,-2)-|node[right,at end]{$-$}(s2);
		\draw[-latex,thick](o3)--++(0,-1)-|node[right,at end]{$-$}(s3);
		\draw[-latex,thick](o2)|-++(-1,2)coordinate(t);
		\node[draw,rectangle,anchor=east] (h1) at(t) {$H_1$};
		\draw[-latex,thick](h1.west)-|node[right,at end]{$-$}(s3);
		\draw[-latex,thick](o1)|-++(5,4)coordinate(t);
		\node[draw,rectangle,anchor=west] (b3) at(t) {$G_3$};
		\draw[-latex,thick](b3.east)-|node[right,at end]{$-$}(s4);
		\onslide<2,8>{\path[flow] (h1.west) -| (s3.center) -- (o2) |- (h1.east);}
		\onslide<3>{\path[flow] (b1.east) -- (o2) --++(0,-2) -| (s2.center) -- (b1.west);}
		\onslide<4>{\path[flow] (b2.east) -- (o3) --++(0,-1) -| (s3.center) -- (b2.west);}
		\onslide<5>{\path[flow] (b1.east) -- (b2) -- (o4) --++(0,-3) -| (s1.center) -- (b1.west);}
		\onslide<6,8>{\path[flow] (b3.east) -| (s4.center) -- (o4) --++(0,-3) -| (s1.center) -- (o1) |- (b3.west);}
		\onslide<7>{\path[flow] (b3.east) -| (s4.center) -- (o3) --++(0,-1) -| (s3.center) -- (o2) --++(0,-2) -| (s2.center) --(o1) |-  (b3.west);}
		\onslide<10>{\path[flow] (s1) -- (b1) -- (b2) -- (o4);}
		\onslide<11>{\path[flow] (s1) -- (o1) |- (b3) -| (s4.center) -- (o4);}
	\end{tikzpicture}
\end{center}
\vskip -1em
\begin{align*}
	\onslide<2->{\Delta&=1-[-H_1}\onslide<3->{-G_1}\onslide<4->{-G_2}\onslide<5->{-G_1G_2}\onslide<6->{+G_3}\onslide<7->{-G_3]}\onslide<8->{-G_3H_1}\\
	\onslide<9->{&=1+H_1+G_1+G_2+G_1G_2-G_3H_1\\}
	\onslide<10->{P_1&=G_1G_2\quad \quad \Delta_1=1\\}
	\onslide<11->{P_2&=-G_3\quad \quad \Delta_2=1+H_1\\}
	\onslide<12->{\Phi&=\frac{G_1G_2-G_3(1+H_1)}{1+H_1+G_1+G_2+G_1G_2-G_3H_1}\\}
	\onslide<13->{&\xlongrightarrow{G_1,G_2,G_3,H_1=\frac{1}{s}}\frac{-1}{s+1}}\onslide<14->{\rightarrow\text{MATLAB求解}}
\end{align*}
\end{frame}

\addtocounter{example}{1}
\begin{frame}{\examplec{求$C(s)/R(s)$：E}{ex:CbyR5}}
\begin{center}
	\begin{tikzpicture}[x=1em, y=1em]
		\draw[-latex,thick](0,0)--node[above,at start]{$R$}++(1em,0)coordinate(o1)--++(1em,0)coordinate(t);
		\node[cross] (s1) at (t)[xshift=0.5em] {};
		\draw[-latex,thick](s1)--++(1,0)coordinate(t);
		\node[draw,rectangle,anchor=west] (b1) at(t) {$G_1$};
		\draw[-latex,thick](b1.east)--++(1,0)coordinate(t);
		\node[cross] (s2) at (t)[xshift=0.5em] {};
		\draw[-latex,thick](s2)--++(1em,0)coordinate(t);
		\node[cross] (s3) at (t)[xshift=0.5em] {};
		\draw[-latex,thick](s3)--++(1em,0)coordinate(t);
		\node[draw,rectangle,anchor=west] (b2) at(t) {$G_2$};
		\draw[-latex,thick](b2.east)--++(1,0)coordinate(o2)--++(1em,0)coordinate(t);
		\node[cross] (s4) at (t)[xshift=0.5em] {};
		\draw[-latex,thick](s4)--++(1,0)coordinate(o3)--++(1em,0)coordinate(t);
		\node[draw,rectangle,anchor=west] (b3) at(t) {$G_3$};
		\draw[-latex,thick](b3.east)--++(1em,0)coordinate(t);
		\node[cross] (s5) at (t)[xshift=0.5em] {};
		\draw[-latex,thick](s5)--++(1,0)coordinate(t);
		\node[draw,rectangle,anchor=west] (b4) at(t) {$G_4$};
		\draw[-latex,thick](b4.east)--++(1,0)coordinate(o4)--node[above,at start]{$C$}++(1em,0)coordinate(t);
		\draw[-latex,thick](o1)|-++(6,3)coordinate(t);
		\node[draw,rectangle,anchor=west] (b6) at (t) {$G_6$};
		\draw[-latex,thick](b6.east)-|node[right,at end]{$-$}(s4);
		\draw[-latex,thick](o1)|-++(2,2)coordinate(t);
		\node[draw,rectangle,anchor=west] (b5) at (t) {$G_5$};
		\draw[-latex,thick](b5.east)-|(s2);
		\draw[-latex,thick](o3)|-++(-3,-2)coordinate(t);
		\node[draw,rectangle,anchor=east] (h2) at (t) {$H_2$};
		\draw[-latex,thick](h2.west)-|(s3);
		\draw[-latex,thick](o4)|-++(-8,-3)coordinate(t);
		\node[draw,rectangle,anchor=east] (h1) at (t) {$H_1$};
		\draw[-latex,thick](h1.west)-|node[right,at end]{$-$}(s1);
		\draw[-latex,thick](o2)--++(0,2em)-|(s5);
		\onslide<2>{\path[flow] (b2) -- (o3) |- (h2) -| (s3.center) -- (b2);}
		\onslide<3>{\path[flow] (b1) -- (b2) -- (b3) -- (b4) -- (o4) |- (h1) -| (s1.center) -- (b1);}
		\onslide<4>{\path[flow] (b1) -- (b2) -- (o2) --++(0,2) -| (s5.center) -- (b4) -- (o4) |- (h1) -| (s1.center) -- (b1);}
		\onslide<6>{\path[flow] (o1) -- (b1) -- (b2) -- (b3) -- (b4) -- (o4);}
		\onslide<7>{\path[flow] (o1) -- (b1) -- (b2) -- (o2) --++ (0,2) -| (s5.center) -- (b4) -- (o4);}
		\onslide<8>{\path[flow] (o1) |- (b5) -| (s2.center) -- (b2) -- (b3) -- (b4) -- (o4);}
		\onslide<9>{\path[flow] (o1) |- (b5) -| (s2.center) -- (b2) -- (o2) --++(0,2) -| (s5) -- (b4) -- (o4);}
		\onslide<10>{\path[flow] (o1) |- (b6) -| (s4) -- (b3) -- (b4) -- (o4);}
		\onslide<11>{\path[flow] (o1) |- (b6) -| (s4) -- (o3) |- (h2) -| (s3) -- (b2) -- (o2) --++(0,2) -| (s5.center) -- (b4) -- (o4);}
	\end{tikzpicture}
\end{center}
\begin{align*}
	\onslide<2->{\Delta&=1-[G_2H_2}\onslide<3->{-G_1G_2G_3G_4H_1}\onslide<4->{-G_1G_2G_4H_1]}\\
	\onslide<5->{&=1-G_2H_2+G_1G_2G_3G_4H_1+G_1G_2G_4H_1\\}
	&\begin{matrix}
		\onslide<6->{P_1=G_1G_2G_3G_4&\Delta_1=1}&\onslide<7->{P_2=G_1G_2G_4&\Delta_2=1\\}
		\onslide<8->{P_3=G_2G_3G_4G_5&\Delta_3=1}&\onslide<9->{P_4=G_2G_4G_5&\Delta_4=1\\}
		\onslide<10->{P_5=-G_3G_4G_6&\Delta_5=1}&\onslide<11->{P_6=-G_6H_2G_2G_4&\Delta_6=1}
	\end{matrix}\\
	\onslide<12->{
	\Phi&=\frac{G_{1-4}+G_{1,2,4}+G_{2-5}+G_{2,4,5}-G_{3,4,6}-G_6H_2G_2G_4}{1-G_2H_2+G_1G_2G_3G_4H_1+G_1G_2G_4H_1}}
\end{align*}
\end{frame}

%\refstepcounter{example}%手动增加计数器计数

\begin{frame}{小结}
\begin{enumerate}
\item 信号流图的组成
\item 信号流图与结构图的关系
\item 结构图$\Rightarrow$信号流图 
\end{enumerate}
\end{frame}

%:Lecture 5: 控制系统的传递函数
\lecture{控制系统的传递函数}{Lecture 5}

\section{控制系统的传递函数}
\subsection{传递函数分类}
\begin{frame}{控制系统的传递函数}
\begin{center}
\tikzstyle{var}=[draw,color=orange,thick,circle,minimum size=0.2em,inner sep=2pt]
\tikzstyle{con}=[draw,color=cyan,thick,-latex]
\tikz{
\draw[-latex,thick](0,0)--node[above,at start]{$R(s)$}++(2em,0)coordinate(t);
\draw[cross] (t)[xshift=0.5em] circle (0.5em)++(0.5em,0)coordinate(t)++(-0.5em,-0.5em)coordinate(c1);
\draw[-latex,thick](t)--node[above,midway]{$E(s)$}++(2em,0)coordinate(t);
\node[draw,rectangle,anchor=west] (b1) at(t) {$G_1(s)$};
\draw[-latex,thick](b1.east)--++(2em,0)coordinate(t);
\draw[cross] (t)[xshift=0.5em] circle (0.5em)++(0.5em,0)coordinate(t)++(-0.5em,0.5em)coordinate(c2);
\draw[-latex,thick](t)--++(2em,0)coordinate(t);
\node[draw,rectangle,anchor=west] (b2) at(t) {$G_2(s)$};
\draw[-latex,thick](b2.east)--++(2em,0)coordinate(o1)--node[above,midway]{$C(s)$}++(2em,0)coordinate(t);
\draw[-latex,thick](o1)|-++(-6em,-3em)coordinate(t);
\node[draw,rectangle,anchor=east] (h) at (t) {$H(s)$};
\draw[-latex,thick](h.west)-|node[right,near end]{$B(s)$}node[right,at end]{$-$}(c1);
\draw[latex-,thick](c2)--node[right,at end]{$N(s)$}++(0,2em)coordinate(t);
} 
\end{center} 
\begin{block}{一、系统的开环传递函数}
定义为把主反馈通道断开，得到的传递函数
$$G_k(s)=\frac{B(s)}{E(s)}=G_1(s)G_2(s)H(s)$$ 
\end{block}
\end{frame}

\begin{frame}{控制系统的传递函数}
\begin{block}{二、输入作用下系统的闭环传递函数}
$$\Phi(s)=\frac{Y(s)}{R(s)}=\frac{G_1(s)G_2(s)}{1+G_1(s)G_2(s)H(s)}$$ 
\end{block}
\begin{block}{三、扰动作用下系统的闭环传递函数}
$$\Phi_N(s)=\frac{Y_N(s)}{N(s)}=\frac{G_2(s)}{1+G_1(s)G_2(s)H(s)}$$ 
\end{block}
\begin{block}{四、系统的总输出}
$$Y(s)=\frac{G_1(s)G_2(s)}{1+G_1(s)G_2(s)H(s)}\cdot R(s)+\frac{G_2(s)}{1+G_1(s)G_2(s)H(s)}\cdot N(s)$$ 
\end{block}
\end{frame}

\begin{frame}{控制系统的传递函数}
\begin{block}{五、误差传递函数}
误差信号$E(s)=R(s)-B(s)$
$$\text{输入作用下的误差传递函数}~~\Phi_e(s)=\frac{E(s)}{R(s)}=\frac{1}{1+G_1(s)G_2(s)H(s)}$$
$$\text{扰动作用下的误差传递函数}~~\Phi_{eN}(s)=\frac{E(s)}{N(s)}=\frac{-G_2(s)H(s)}{1+G_1(s)G_2(s)H(s)}$$
\end{block}
\begin{block}{六、系统的总误差}
$$E(s)=\frac{R(s)-G_2(s)H(s)N(s)}{1+G_1(s)G_2(s)H(s)}$$ 
\end{block}
\end{frame}

\subsection{求传递函数实例}
\begin{frame}{\example{求$C(s)/R(s)$，$C(s)/N(s)$}{ex:CbyR2}}
\begin{center}
\tikzstyle{var}=[draw,color=orange,thick,circle,minimum size=0.2em,inner sep=2pt]
\tikzstyle{con}=[draw,color=cyan,thick,-latex]
\begin{tikzpicture}
\draw[-latex,thick](0,0)--node[above,at start]{$R(s)$}++(1em,0)coordinate(t);
\draw[cross] (t)[xshift=0.5em] circle (0.5em)++(0.5em,0)coordinate(t)++(-0.5em,-0.5em)coordinate(c1);
\draw[-latex,thick](t)--++(1em,0)coordinate(t);
\node[draw,rectangle,anchor=west] (b1) at(t) {$G_1$};
\draw[-latex,thick](b1.east)--++(1em,0)coordinate(o1)--++(1em,0)coordinate(t);
\draw[cross] (t)[xshift=0.5em] circle (0.5em)++(0.5em,0)coordinate(t)++(-0.5em,-0.5em)coordinate(c2);
\draw[-latex,thick](t)--++(1em,0)coordinate(t);
\node[draw,rectangle,anchor=west] (b2) at(t) {$G_2$};
\draw[-latex,thick](b2.east)--++(1em,0)coordinate(o2)--++(1em,0)coordinate(t);
\draw[cross] (t)[xshift=0.5em] circle (0.5em)++(0.5em,0)coordinate(t)++(-0.5em,-0.5em)coordinate(c3);
\draw[-latex,thick](t)--++(1em,0)coordinate(o3)--node[above,at end]{$C(s)$}++(1em,0)coordinate(t);
\draw[-latex,thick](o2)|-++(-2em,-2em)coordinate(t);
\node[draw,rectangle,anchor=east] (h) at (t) {$H$};
\draw[-latex,thick](h.west)-|node[right,at end]{$-$}(c2);
\draw[-latex,thick](o1)|-++(4em,-3em)coordinate(t);
\node[draw,rectangle,anchor=west] (b3) at (t) {$G_3$};
\draw[-latex,thick](b3.east)-|(c3);
\draw[-latex,thick](o3)--++(0,-4em)-|node[right,at end]{$-$}(c1);
\draw[-latex,thick]($(c3)+(0,3em)$)|-++(-5em,-1em)coordinate(t);
\node[draw,rectangle,anchor=east] (b4) at (t) {$G_4$};
\draw[-latex,thick](b4.west)-|($(c1)+(0,1em)$);
\draw[-latex,thick]($(c3)+(0,3em)$)--node[right,at start]{$N(s)$}node[left,very near end]{$-$}($(c3)+(0,1em)$);
\end{tikzpicture}
\end{center}
\begin{align*}
\Delta&=1-[-G_2H-G_1G_2-G_1G_3]+G_1G_2G_3H\\
&=1+G_2H+G_1G_2+G_1G_3+G_1G_2G_3H
\end{align*}
\begin{equation*}\left.
\begin{matrix}
P_1=G_1G_2&\Delta_1=1\\
P_2=G_1G_3&\Delta_2=1+G_2H
\end{matrix}\right\}
\Phi=\frac{G_1G_2+G_1G_3(1+G_2H)}{1+G_2H+G_1G_2+G_1G_3+G_{1-3}H}
\end{equation*}
\begin{equation*}\left.
\begin{matrix}
P_{N1}=-1&\Delta_{N1}=1+G_2H\\
P_{N2}=G_{4,1,2}&\Delta_{N2}=1\\
P_{N3}=G_{4,1,3}&\Delta_{N3}=1+G_2H
\end{matrix}\right\}
\Phi_N=\frac{(-1+G_{1,3,4})(1+G_2H)+G_{1,2,4}}{1+G_2H+G_{1,2}+G_{1,3}+G_{1-3}H}
\end{equation*}
\end{frame}

\begin{frame}{\example{梅逊公式求$C(s)$}{ex:mason1}}
\begin{center}
\begin{tikzpicture}
\draw[-latex,thick](0,0)--node[above,at start]{$R(s)$}++(1em,0)coordinate(o1)--++(1em,0)coordinate(t);
\draw[cross] (t)[xshift=0.5em] circle (0.5em)++(0.5em,0)coordinate(t)++(-0.5em,-0.5em)coordinate(c1);
\draw[-latex,thick](t)--node[above,midway]{$E(s)$}++(2em,0)coordinate(t);
\node[draw,rectangle,anchor=west] (b1) at(t) {$G_1(s)$};
\draw[-latex,thick](b1.east)--++(1em,0)coordinate(o2)--++(1em,0)coordinate(t);
\draw[cross] (t)[xshift=0.5em] circle (0.5em)++(0.5em,0)coordinate(t)++(-0.5em,-0.5em)coordinate(c2)++(0,1em)coordinate(c3);
\draw[-latex,thick](t)--++(1em,0)coordinate(t);
\draw[cross] (t)[xshift=0.5em] circle (0.5em)++(0.5em,0)coordinate(t)++(-0.5em,0.5em)coordinate(c4);
\draw[-latex,thick](t)--++(1em,0)coordinate(t);
\node[draw,rectangle,anchor=west] (b2) at(t) {$G_2(s)$};
\draw[-latex,thick](b2.east)--++(1em,0)coordinate(o3)--node[above,at end]{$C(s)$}++(1em,0)coordinate(t);
\draw[cross] ($(c1)+(0,-2em)$)circle (0.5em)++(0,-0.5em)coordinate(c5)++(0.5em,0.5em)coordinate(c6);
\draw[-latex,thick](o2)|-++(-1em,-2.5em)coordinate(t);
\node[draw,rectangle,anchor=east] (h1) at (t) {$H_1(s)$};
\draw[-latex,thick](h1.west)--node[above,near end]{$-$}(c6);
\draw[-latex,thick]($(c5)+(0,1em)$)--node[right,at end]{$-$}(c1);
\draw[-latex,thick](o3)|-++(-2em,-2.5em)coordinate(t);
\node[draw,rectangle,anchor=east] (h2) at (t) {$H_2(s)$};
\draw[-latex,thick](h2.west)-|node[right,at end]{$-$}(c2);
\draw[-latex,thick](o3)|-++(-6em,-4em)coordinate(t);
\node[draw,rectangle,anchor=east] (h3) at (t) {$H_3(s)$};
\draw[-latex,thick](h3.west)-|(c5);
\draw[-latex,thick](o1)|-++(3em,2.5em)coordinate(t);
\node[draw,rectangle,anchor=west] (b3) at (t) {$G_3(s)$};
\draw[-latex,thick](b3.east)-|(c3);
\draw[-latex,thick]($(c4)+(0,3em)$)--node[right,at start]{$N(s)$}(c4);
\end{tikzpicture}
\end{center}
\begin{equation*}
\begin{matrix}
L_1=G_1H_1&L_2=-G_2H_2\\
L_3=-G_1G_2H_3&L_1L_2=(G_1H_1)(-G_2H_2)
\end{matrix}
\end{equation*}
$$C(s)=\frac{\alert{R(s)[}G_3G_2(1-G_1H_1)+G_1G_2\alert{]}+G_2(1-G_1H_1)\alert{N(s)}}{1-G_1H_1+G_2H_2+G_1G_2H_3-G_1H_1G_2H_2}$$
\end{frame}

\begin{frame}{\example{梅逊公式求$E(s)$}{ex:mason2}}
\begin{center}
\begin{tikzpicture}
\draw[-latex,thick](0,0)--node[above,at start]{$R(s)$}++(1em,0)coordinate(o1)--++(1em,0)coordinate(t);
\draw[cross] (t)[xshift=0.5em] circle (0.5em)++(0.5em,0)coordinate(t)++(-0.5em,-0.5em)coordinate(c1);
\draw[-latex,thick,color=red]($(t)+(2em,0)$)--($(t)+(4em,1em)$);
\draw[-latex,thick](t)--node[above,near start]{$E(s)$}++(4em,0)coordinate(t);
\node[draw,rectangle,anchor=west] (b1) at(t) {$G_1(s)$};
\draw[-latex,thick](b1.east)--++(1em,0)coordinate(o2)--++(1em,0)coordinate(t);
\draw[cross] (t)[xshift=0.5em] circle (0.5em)++(0.5em,0)coordinate(t)++(-0.5em,-0.5em)coordinate(c2)++(0,1em)coordinate(c3);
\draw[-latex,thick](t)--++(1em,0)coordinate(t);
\draw[cross] (t)[xshift=0.5em] circle (0.5em)++(0.5em,0)coordinate(t)++(-0.5em,0.5em)coordinate(c4);
\draw[-latex,thick](t)--++(1em,0)coordinate(t);
\node[draw,rectangle,anchor=west] (b2) at(t) {$G_2(s)$};
\draw[-latex,thick](b2.east)--++(1em,0)coordinate(o3)--node[above,at end]{$C(s)$}++(1em,0)coordinate(t);
\draw[cross] ($(c1)+(0,-2em)$)circle (0.5em)++(0,-0.5em)coordinate(c5)++(0.5em,0.5em)coordinate(c6);
\draw[-latex,thick](o2)|-++(-1em,-2.5em)coordinate(t);
\node[draw,rectangle,anchor=east] (h1) at (t) {$H_1(s)$};
\draw[-latex,thick](h1.west)--node[above,near end]{$-$}(c6);
\draw[-latex,thick]($(c5)+(0,1em)$)--node[right,at end]{$-$}(c1);
\draw[-latex,thick](o3)|-++(-2em,-2.5em)coordinate(t);
\node[draw,rectangle,anchor=east] (h2) at (t) {$H_2(s)$};
\draw[-latex,thick](h2.west)-|node[right,at end]{$-$}(c2);
\draw[-latex,thick](o3)|-++(-6em,-4em)coordinate(t);
\node[draw,rectangle,anchor=east] (h3) at (t) {$H_3(s)$};
\draw[-latex,thick](h3.west)-|(c5);
\draw[-latex,thick](o1)|-++(3em,2.5em)coordinate(t);
\node[draw,rectangle,anchor=west] (b3) at (t) {$G_3(s)$};
\draw[-latex,thick](b3.east)-|(c3);
\draw[-latex,thick]($(c4)+(0,3em)$)--node[right,at start]{$N(s)$}(c4);
\end{tikzpicture}
\end{center}
\begin{equation*}
\begin{matrix}
P_1=1&\Delta_1=1+G_2H_2\\
P_2=-G_3G_2H_3&\Delta_2=1\\
\alert{P_1=-G_2H_3}&\alert{\Delta_1=1}
\end{matrix}
\end{equation*}
$$E(s)=\frac{\alert{R(s)[}(1+G_2H_2)+(-G_3G_2H_3)\alert{]}+(-G_2H_3)\alert{N(s)}}{1-G_1H_1+G_2H_2+G_1G_2H_3-G_1H_1G_2H_2}$$
\end{frame}

\begin{frame}{本章小结}
\begin{itemize}
\item 掌握建立微分方程的方法
\item 掌握拉氏变换求解微分方程的方法
\item 牢固掌握系统传递函数的定义
\item 能熟练地进行动态结构图等效变换
\item 能熟练运用梅逊公式求取系统传递函数
\item 了解控制系统中各种传递函数的定义 
\end{itemize}
\end{frame}

\end{document}